\documentclass[11pt,a4paper]{article}

% ==================== TEMEL PAKETLER (pdflatex) ====================
\usepackage[T1]{fontenc}
\usepackage[utf8]{inputenc}
\usepackage{lmodern}
\usepackage[a4paper,margin=2cm,top=2.4cm,bottom=2.4cm]{geometry}
\usepackage{xcolor}
\usepackage{colortbl}
\usepackage{array}
\usepackage{booktabs}
\usepackage{longtable}
\usepackage{tabularx}
\usepackage{listings}
\usepackage{fancyhdr}
\usepackage{amsmath}
\usepackage{amssymb}
\usepackage{hyperref}

% ==================== RENK PALETİ ====================
\definecolor{anaMavi}{HTML}{132C4A}
\definecolor{vurguMavi}{HTML}{2E6DA4}
\definecolor{acikMavi}{HTML}{EAF2F9}
\definecolor{turuncu}{HTML}{D9701E}
\definecolor{yesil}{HTML}{2E7D32}
\definecolor{acikYesil}{HTML}{E8F5E9}
\definecolor{kirmizi}{HTML}{C0392B}
\definecolor{mor}{HTML}{6C3483}
\definecolor{acikMor}{HTML}{F3E9F7}
\definecolor{acikGri}{HTML}{F4F5F7}
\definecolor{ortaGri}{HTML}{6B7280}
\definecolor{koyuGri}{HTML}{2D2D2D}
\definecolor{kodArka}{HTML}{FBFBFA}

% ==================== HYPERREF ====================
\hypersetup{
  colorlinks=true,
  linkcolor={[HTML]{2E6DA4}},
  urlcolor={[HTML]{2E6DA4}},
  pdftitle={C++ Programlama Dili Kapsamlı Rehber},
  pdfauthor={Modern C++ Rehberi},
  bookmarksnumbered=true
}

% ==================== BAŞLIK STİLLERİ (titlesec olmadan) ====================
\setcounter{secnumdepth}{0}
\newcounter{secc}
\newcounter{subsecc}[secc]
\newcounter{subsubsecc}[subsecc]
\newcommand{\gsec}[1]{\par\phantomsection\stepcounter{secc}%
  \setcounter{subsecc}{0}%
  \vspace{18pt}{\color{anaMavi}\Large\bfseries\thesecc.\hspace{0.5em}#1}\par
  \vspace{2pt}{\color{turuncu}\rule{\linewidth}{1.6pt}}\par\vspace{8pt}%
  \addcontentsline{toc}{section}{\protect\numberline{\thesecc}#1}\markboth{#1}{#1}}
\newcommand{\gsub}[1]{\par\phantomsection\stepcounter{subsecc}%
  \setcounter{subsubsecc}{0}%
  \vspace{11pt}{\color{vurguMavi}\large\bfseries\thesecc.\thesubsecc\hspace{0.5em}#1}\par\vspace{4pt}%
  \addcontentsline{toc}{subsection}{\protect\numberline{\thesecc.\thesubsecc}#1}}
\newcommand{\gsubsub}[1]{\par\phantomsection\stepcounter{subsubsecc}%
  \vspace{7pt}{\color{koyuGri}\normalsize\bfseries #1}\par\vspace{2pt}}

% Kısım (Part) başlığı
\newcommand{\gpart}[2]{\clearpage
  \begin{center}\vspace*{1.5cm}
  {\color{turuncu}\rule{0.7\linewidth}{1.5pt}}\\[10pt]
  {\color{ortaGri}\Large KISIM #1}\\[8pt]
  {\color{anaMavi}\Huge\bfseries #2}\\[10pt]
  {\color{turuncu}\rule{0.7\linewidth}{1.5pt}}
  \end{center}\vspace{1cm}}

% ==================== BAŞLIK / ALTBİLGİ ====================
\setlength{\headheight}{14pt}
\pagestyle{fancy}
\fancyhf{}
\renewcommand{\headrulewidth}{0.4pt}
\renewcommand{\footrulewidth}{0.4pt}
\fancyhead[L]{\small\color{ortaGri}C++ Programlama Rehberi}
\fancyhead[R]{\small\color{ortaGri}\leftmark}
\fancyfoot[C]{\small\color{ortaGri}\thepage}

% ==================== KOD LİSTELEME (C++) ====================
\lstdefinestyle{cppstyle}{
  language=C++,
  basicstyle=\ttfamily\small\color{koyuGri},
  keywordstyle=\color{vurguMavi}\bfseries,
  commentstyle=\color{yesil}\itshape,
  stringstyle=\color{turuncu},
  numberstyle=\tiny\color{ortaGri},
  numbers=left,
  numbersep=8pt,
  backgroundcolor=\color{kodArka},
  frame=single,
  framesep=6pt,
  rulecolor=\color{ortaGri!40},
  breaklines=true,
  showstringspaces=false,
  tabsize=4,
  columns=fullflexible,
  keepspaces=true,
  extendedchars=true,
  literate={ı}{{\i}}1{İ}{{\.I}}1{ş}{{\c s}}1{Ş}{{\c S}}1{ğ}{{\u g}}1{Ğ}{{\u G}}1{ç}{{\c c}}1{Ç}{{\c C}}1{ö}{{\"o}}1{Ö}{{\"O}}1{ü}{{\"u}}1{Ü}{{\"U}}1,
  morekeywords={size_t,uint8_t,uint16_t,uint32_t,uint64_t,int8_t,int16_t,
    int32_t,int64_t,nullptr,constexpr,consteval,noexcept,override,final,auto,
    vector,string,pair,map,set,unordered_map,unordered_set,priority_queue,
    queue,stack,deque,array,tuple,optional,shared_ptr,unique_ptr,weak_ptr,
    template,typename,class,public,private,protected,virtual,namespace,
    using,std,move,forward,make_pair,make_shared,make_unique,decltype,
    default,delete,explicit,mutable,friend,operator,noexcept}
}
\lstset{style=cppstyle}

% Kabuk/terminal stili
\lstdefinestyle{shstyle}{
  basicstyle=\ttfamily\small\color{koyuGri},
  backgroundcolor=\color{koyuGri!6},
  frame=single,framesep=6pt,rulecolor=\color{ortaGri!40},
  breaklines=true,showstringspaces=false,tabsize=4,columns=fullflexible,
  keepspaces=true,
  literate={ı}{{\i}}1{İ}{{\.I}}1{ş}{{\c s}}1{Ş}{{\c S}}1{ğ}{{\u g}}1{Ğ}{{\u G}}1{ç}{{\c c}}1{Ç}{{\c C}}1{ö}{{\"o}}1{Ö}{{\"O}}1{ü}{{\"u}}1{Ü}{{\"U}}1,
}

% ==================== ÖZEL KUTULAR (tcolorbox olmadan) ====================
\newcommand{\callout}[4]{\par\smallskip\noindent\begingroup
  \setlength{\fboxsep}{8pt}\setlength{\fboxrule}{0.8pt}%
  \fcolorbox{#1}{#2}{\parbox{\dimexpr\linewidth-2\fboxsep-2\fboxrule\relax}{%
    {\bfseries\color{#1}#3}\par\smallskip #4}}\endgroup\par\smallskip}
\newcommand{\bilgi}[1]{\callout{vurguMavi}{acikMavi}{Bilgi}{#1}}
\newcommand{\ipucu}[1]{\callout{yesil}{acikYesil}{İpucu}{#1}}
\newcommand{\uyari}[1]{\callout{kirmizi}{kirmizi!7}{Dikkat}{#1}}
\newcommand{\ornek}[2][Örnek]{\callout{ortaGri!90!black}{acikGri}{#1}{#2}}
\newcommand{\notk}[2][Not]{\callout{mor}{acikMor}{#1}{#2}}

% Soru / Çözüm sistemi (alıştırmalar için)
\newcounter{soruc}
\newcommand{\soru}[1]{\par\smallskip\refstepcounter{soruc}%
  \noindent\begingroup\setlength{\fboxsep}{8pt}\setlength{\fboxrule}{0.8pt}%
  \fcolorbox{turuncu}{turuncu!6}{\parbox{\dimexpr\linewidth-2\fboxsep-2\fboxrule\relax}{%
  {\bfseries\color{turuncu}Soru \thesoruc}\par\smallskip #1}}\endgroup\par\smallskip}
\newcommand{\cozum}{\par\noindent{\bfseries\color{yesil}Çözüm.}\ }

% Yardımcı komutlar
\newcommand{\code}[1]{\texttt{\color{koyuGri}#1}}
\newcommand{\bigO}[1]{\ensuremath{\mathcal{O}(#1)}}
\let\svtimes\times   \renewcommand{\times}{\ensuremath{\svtimes}}
\let\svcdot\cdot     \renewcommand{\cdot}{\ensuremath{\svcdot}}
\let\svto\to         \renewcommand{\to}{\ensuremath{\svto}}
\let\svle\le         \renewcommand{\le}{\ensuremath{\svle}}
\let\svge\ge         \renewcommand{\ge}{\ensuremath{\svge}}
\renewcommand{\arraystretch}{1.35}

% ==================== BAŞLANGIÇ ====================
\begin{document}

% ---------- KAPAK ----------
\begin{titlepage}
\centering
\vspace*{1.4cm}
{\color{turuncu}\rule{\textwidth}{2pt}}\\[8pt]
{\color{anaMavi}\Huge\bfseries C++ Programlama Dili}\\[10pt]
{\color{vurguMavi}\LARGE Kapsamlı Uygulamalı Rehber}\\[8pt]
{\color{ortaGri}\large Temellerden Modern C++'a: Nesne Yönelimi, Bellek Yönetimi, Şablonlar ve STL}\\[6pt]
{\color{turuncu}\rule{\textwidth}{2pt}}\\[1.4cm]

\begin{center}
\fcolorbox{vurguMavi}{acikMavi}{\parbox{0.86\textwidth}{\centering\large
Bu belge, C++'ı temel sözdiziminden başlayarak nesne yönelimli programlamaya,
bellek ve kaynak yönetiminin inceliklerine, şablonlara, STL'e ve modern
C++17/20 özelliklerine kadar \textbf{derinlemesine ve çalışan kod
örnekleriyle} anlatır. Her sınıfın kalbi olan \textbf{özel üye fonksiyonlar}
(Beşin Kuralı), akıllı işaretçiler, istisna güvenliği ve operatör aşırı
yükleme gibi kritik konular ayrıntılı biçimde ele alınır.\vspace{4pt}}}
\end{center}

\vspace{0.7cm}
\begin{center}
\fcolorbox{ortaGri!50}{acikGri}{\parbox{0.86\textwidth}{\small
\textbf{K1 (Temeller):} program yapısı \textbullet\ tipler \& \code{auto}
\textbullet\ kontrol akışı \textbullet\ fonksiyonlar.\quad
\textbf{K2 (OOP):} sınıflar \textbullet\ kalıtım \textbullet\ polimorfizm
\& sanal fonksiyonlar.\vspace{3pt}

\textbf{K3 (Bellek \& İşaretçiler):} yığın/öbek \textbullet\ \code{new}/\code{delete}
\textbullet\ işaretçi vs referans \textbullet\ \code{const} ile işaretçiler.\quad
\textbf{K4 (RAII \& Akıllı İşaretçiler):} \code{unique\_ptr} \textbullet\
\code{shared\_ptr} \textbullet\ \code{weak\_ptr}.\vspace{3pt}

\textbf{K5 (Beşin Kuralı):} lvalue/rvalue \& \code{std::move} \textbullet\
yıkıcı \textbullet\ kopya/taşıma yapıcı \& atama \textbullet\ copy-and-swap.\quad
\textbf{K6 (Şablonlar):} fonksiyon/sınıf şablonları \textbullet\ variadic
\textbullet\ özelleştirme.\vspace{3pt}

\textbf{K7 (STL):} konteynerler \textbullet\ iteratörler \textbullet\
algoritmalar.\quad
\textbf{K8 (Lambda):} lambda ifadeleri \textbullet\ \code{std::function}.\quad
\textbf{K9 (İstisnalar):} \code{try}/\code{catch} \textbullet\ istisna güvenliği.\vspace{3pt}

\textbf{K10 (Anahtar Kelimeler):} \code{explicit} \textbullet\ \code{constexpr}
\textbullet\ \code{noexcept} \textbullet\ \code{static}/\code{mutable}/\code{friend}.\quad
\textbf{K11 (Operatörler):} aritmetik \textbullet\ \code{<=>} \textbullet\
\code{[]}, \code{()} \textbullet\ dönüşüm.\quad
\textbf{K12 (Modern C++):} yapısal bağlama \textbullet\ concepts \& ranges.\vspace{4pt}}}
\end{center}

\vfill
{\color{ortaGri}\large Hazırlanma Tarihi: 19 Temmuz 2026}\\[4pt]
{\color{ortaGri} C++17 / C++20 \textbullet\ g++ / clang++ \textbullet\ -Wall -Wextra}
\vspace*{0.6cm}
\end{titlepage}

% ---------- İÇİNDEKİLER ----------
\tableofcontents
\thispagestyle{empty}
\clearpage

% #####################################################################
\gpart{1}{C++ Temelleri}
% #####################################################################

% =====================================================================
\gsec{Bir C++ Programının Yapısı}

C++, C dilinin üzerine inşa edilmiş, çok paradigmalı (yordamsal, nesne
yönelimli, jenerik ve fonksiyonel) bir sistem programlama dilidir. Derlenen
(compiled) bir dildir: kaynak kod, doğrudan makine koduna çevrilir; bu da
onu son derece hızlı kılar. En küçük çalışan programla başlayalım:

\begin{lstlisting}
#include <iostream>    // std::cout icin standart giris/cikis akisi

// main: her programin baslangic noktasidir; isletim sistemi burayi cagirir
int main() {
    std::cout << "Merhaba, Dunya!" << "\n";   // ekrana yaz
    return 0;          // 0 = basarili cikis kodu
}
\end{lstlisting}

\begin{lstlisting}[style=shstyle]
# Derleme ve calistirma (GCC)
g++ -std=c++17 -Wall -Wextra merhaba.cpp -o merhaba
./merhaba
# Cikti: Merhaba, Dunya!
\end{lstlisting}

\bilgi{Derleme dört aşamadan geçer: \textbf{ön işleme} (\code{\#include},
makrolar), \textbf{derleme} (kaynak $\to$ assembly), \textbf{çevirme}
(assembly $\to$ nesne dosyası \code{.o}) ve \textbf{bağlama} (linking:
nesne dosyaları + kütüphaneler $\to$ çalıştırılabilir). \code{-Wall -Wextra}
bayraklarını her zaman açın; potansiyel hataları derleme aşamasında yakalar.}

\gsub{Başlık Dosyaları ve Ad Alanları (Namespaces)}

Kod, \emph{bildirimlerin} bulunduğu başlık dosyaları (\code{.hpp}) ve
\emph{tanımların} bulunduğu kaynak dosyalar (\code{.cpp}) olarak
örgütlenir. \code{namespace}, isim çakışmalarını önlemek için sembolleri
mantıksal gruplara ayırır.

\begin{lstlisting}
#include <iostream>

// Kendi ad alanimiz -- isim cakismalarini onler
namespace math {
    const double PI = 3.14159265358979;
    double square(double x) { return x * x; }

    namespace geometry {          // ic ice ad alani
        double circle_area(double r) { return PI * square(r); }
    }
}

int main() {
    // Tam nitelenmis isim (fully qualified)
    std::cout << math::geometry::circle_area(2.0) << "\n";

    // 'using' ile bir ismi kapsama getir
    using math::square;
    std::cout << square(5) << "\n";       // 25

    // DIKKAT: 'using namespace std;' baslik dosyalarinda ASLA kullanilmaz
    // (tum std'yi getirir, cakismalara yol acar). .cpp icinde olc

    return 0;
}
\end{lstlisting}

\uyari{Başlık dosyalarında (\code{.hpp}) asla \code{using namespace std;}
yazmayın. Bu satır, o başlığı dahil eden \emph{her} dosyaya tüm \code{std}
ad alanını enjekte eder ve gizli isim çakışmalarına yol açar. Kaynak
dosyalarda (\code{.cpp}) ölçülü kullanılabilir; en güvenlisi her zaman
\code{std::} önekini açıkça yazmaktır.}

% =====================================================================
\gsec{Temel Tipler, Değişkenler ve auto}

C++ \emph{statik tipli} bir dildir: her değişkenin tipi derleme zamanında
bilinir ve sabittir. Temel (yerleşik) tipler bellekte belirli boyutlar
kaplar ve belirli değer aralıklarına sahiptir.

\begin{center}
\begin{tabular}{@{}lll@{}}
\toprule
\textbf{Tip} & \textbf{Tipik Boyut} & \textbf{Açıklama} \\
\midrule
\code{bool} & 1 bayt & \code{true} / \code{false} \\
\code{char} & 1 bayt & Karakter / küçük tam sayı \\
\code{int} & 4 bayt & Tam sayı ($\approx \pm 2.1$ milyar) \\
\code{long long} & 8 bayt & Büyük tam sayı \\
\code{float} & 4 bayt & Tek hassasiyetli ondalık \\
\code{double} & 8 bayt & Çift hassasiyetli ondalık \\
\code{std::size\_t} & 8 bayt & İşaretsiz boyut/indeks tipi \\
\bottomrule
\end{tabular}
\end{center}

\begin{lstlisting}
#include <iostream>
#include <string>
#include <cstdint>

int main() {
    // Modern C++: kume parantezli baslatma (uniform initialization)
    int    age{25};            // brace-init: daraltmaya (narrowing) izin vermez
    double pi{3.14159};
    bool   active{true};
    char   letter{'A'};
    std::string name{"Ada"};

    // int hata{3.14};        // HATA: double -> int daraltma engellenir (guvenli!)

    // 'auto': tipi derleyici initializer'dan CIKARIR
    auto number = 42;            // int
    auto ratio = 3.5;          // double
    auto text = std::string{"otomatik"};  // std::string
    auto& ref = age;          // int& (referans icin & yaz)

    // Tamsayi literalleri ve ayirici
    auto big = 1'000'000;   // ' okunabilirlik ayiricisi (C++14)
    std::uint64_t mask = 0xFF00;   // onaltilik (hex)

    std::cout << name << ", " << age << " (auto sayi=" << number << ")\n";

    // const ve constexpr
    const int CONSTANT = 100;             // calisma zamani sabiti
    constexpr int COMPILE_CONSTANT = 50; // derleme zamani sabiti
    std::cout << CONSTANT + COMPILE_CONSTANT << "\n";
    return 0;
}
\end{lstlisting}

\ipucu{\code{auto}, kodu kısaltır ve tip tekrarını önler; özellikle uzun
iteratör tiplerinde (\code{std::vector<int>::iterator} yerine \code{auto})
çok değerlidir. Ancak niyeti belirsizleştirdiği yerlerde açık tip yazmak
daha okunur olabilir. Kume parantezli başlatma (\code{\{\}}) ise ``daraltma''
(narrowing) dönüşümlerini derleme zamanında yakaladığı için en güvenli
başlatma biçimidir.}

% =====================================================================
\gsec{Kontrol Akışı}

\begin{lstlisting}
#include <iostream>
#include <vector>
#include <string>

int main() {
    int score = 75;

    // if / else if / else
    if (score >= 90)      std::cout << "AA\n";
    else if (score >= 70) std::cout << "BB\n";
    else                 std::cout << "Kaldi\n";

    // switch: tam sayi/enum degerlerine dallan
    int day = 3;
    switch (day) {
        case 1: std::cout << "Pazartesi\n"; break;
        case 3: std::cout << "Carsamba\n";  break;
        default: std::cout << "Diger\n";    break;
    }

    // Klasik for
    for (int i = 0; i < 5; ++i) std::cout << i << " ";
    std::cout << "\n";

    // while ve do-while
    int n = 3;
    while (n > 0) { std::cout << n << " "; --n; }
    std::cout << "\n";

    // RANGE-BASED FOR (C++11): en cok tercih edilen dongu bicimi
    std::vector<std::string> colors{"kirmizi", "yesil", "mavi"};
    for (const auto& color : colors)      // kopyalamadan gez (const ref)
        std::cout << color << " ";
    std::cout << "\n";

    // Indeksle gezmek gerekiyorsa yapısal bağlama ile (C++17)
    for (std::size_t i = 0; i < colors.size(); ++i)
        std::cout << i << ":" << colors[i] << " ";
    std::cout << "\n";
    return 0;
}
\end{lstlisting}

\bilgi{\textbf{Range-based for} (\code{for (auto\& x : kap)}) modern C++'ın
tercih edilen döngüsüdür: indeks hatası (off-by-one) riskini ortadan
kaldırır ve niyeti nettir. Elemanları değiştirmeyecekseniz \code{const auto\&},
değiştirecekseniz \code{auto\&}, ucuz kopyalanabilir küçük tipler için
\code{auto} kullanın.}

% =====================================================================
\gsec{Fonksiyonlar}

Fonksiyonlar, kodu yeniden kullanılabilir birimlere böler. C++'ta
parametreler değere, referansa veya sabit referansa göre geçirilebilir ---
bu seçim hem performansı hem anlamı belirler.

\begin{lstlisting}
#include <iostream>
#include <string>
#include <vector>

// 1) DEGERE gore: kopya alir; kucuk tipler icin uygun
int square(int x) { return x * x; }

// 2) REFERANSA gore: cagiranin degiskenini DEGISTIRIR
void double_value(int& x) { x *= 2; }

// 3) SABIT REFERANSA gore: buyuk nesneyi KOPYALAMADAN, degistirmeden gecir
double average(const std::vector<double>& values) {
    double sum = 0;
    for (double v : values) sum += v;
    return values.empty() ? 0.0 : sum / values.size();
}

// 4) VARSAYILAN argumanlar
void greet(const std::string& name, const std::string& title = "Sn.") {
    std::cout << title << " " << name << ", merhaba!\n";
}

// 5) ASIRI YUKLEME (overloading): ayni isim, farkli parametre listeleri
int add(int a, int b)          { return a + b; }
double add(double a, double b) { return a + b; }

// 6) Sondaki donus tipi (trailing return type) -- karmasik tiplerde
auto multiply(int a, int b) -> int { return a * b; }

int main() {
    std::cout << square(6) << "\n";                 // 36

    int value = 10;
    double_value(value);                          // referans -> deger degisir
    std::cout << value << "\n";                    // 20

    std::cout << average({1.0, 2.0, 3.0, 4.0}) << "\n";  // 2.5

    greet("Ada");                                  // varsayilan unvan
    greet("Turing", "Prof.");

    std::cout << add(2, 3) << " " << add(2.5, 3.5) << "\n";  // 5  6
    std::cout << multiply(4, 5) << "\n";           // 20
    return 0;
}
\end{lstlisting}

\ipucu{\textbf{Parametre geçirme kuralı:} Küçük, ucuz kopyalanan tipler
(\code{int}, \code{double}, işaretçi) için \emph{değere} geçirin. Büyük
nesneleri (\code{std::string}, \code{std::vector}, sınıflar) yalnızca
okuyacaksanız \code{const T\&} ile geçirin --- kopya maliyetinden kaçınır.
Fonksiyonun çağıranın nesnesini \emph{değiştirmesi} gerekiyorsa \code{T\&}
(referans) kullanın.}

% #####################################################################
\gpart{2}{Nesne Yönelimli Programlama}
% #####################################################################

% =====================================================================
\gsec{Sınıflar ve Nesneler}

Sınıf (class), veriyi (üye değişkenler) ve o veri üzerinde çalışan
davranışı (üye fonksiyonlar) tek bir birimde toplayan kullanıcı tanımlı
bir tiptir. Bu birleştirmeye \textbf{kapsülleme} (encapsulation) denir.
Erişim belirleyiciler (\code{private}, \code{public}, \code{protected})
verinin dışarıdan nasıl görüneceğini kontrol eder.

\begin{lstlisting}
#include <iostream>
#include <string>
#include <stdexcept>

class BankAccount {
private:                          // disaridan ERISILEMEZ (kapsulleme)
    std::string owner_;
    double balance_;

public:                           // genel arayuz (interface)
    // Yapici (constructor): nesne olustugunda uyeleri baslatir
    BankAccount(std::string owner, double initial = 0.0)
        : owner_(std::move(owner)), balance_(initial) {
        if (initial < 0) throw std::invalid_argument("Negatif bakiye");
    }

    // Uye fonksiyonlar (metotlar) -- veriyi kontrollu degistirir
    void deposit(double amount) {
        if (amount <= 0) throw std::invalid_argument("Gecersiz miktar");
        balance_ += amount;
    }

    void withdraw(double amount) {
        if (amount > balance_) throw std::runtime_error("Yetersiz bakiye");
        balance_ -= amount;
    }

    // const uye fonksiyon: nesneyi DEGISTIRMEZ (yalnizca okur)
    double balance() const { return balance_; }
    const std::string& owner() const { return owner_; }
};

int main() {
    BankAccount account("Ada Lovelace", 1000);
    account.deposit(500);
    account.withdraw(200);
    std::cout << account.owner() << " bakiye: " << account.balance() << "\n"; // 1300

    try {
        account.withdraw(999999);        // yetersiz bakiye -> istisna
    } catch (const std::exception& e) {
        std::cout << "Hata: " << e.what() << "\n";
    }
    return 0;
}
\end{lstlisting}

\bilgi{\textbf{Kapsülleme ilkesi:} Üye değişkenleri \code{private} yapın ve
dışarıya yalnızca kontrollü bir \code{public} arayüz sunun. Böylece sınıfın
iç durumu her zaman tutarlı (invariant) kalır --- yukarıda bakiye asla
negatif olamaz, çünkü tek değiştirme yolu \code{yatir}/\code{cek}
metotlarından geçer ve onlar doğrulama yapar. \code{struct} ile \code{class}
teknik olarak yalnızca varsayılan erişim belirleyicisinde farklıdır
(\code{struct}: public, \code{class}: private).}

% =====================================================================
\gsec{Kalıtım (Inheritance)}

Kalıtım, mevcut bir sınıftan (\emph{taban}, base) yeni bir sınıf
(\emph{türetilmiş}, derived) oluşturmayı sağlar. Türetilmiş sınıf, tabanın
üyelerini devralır ve genişletir. Bu, ``bir-türüdür'' (is-a) ilişkisini
modeller: bir \code{Kopek}, bir \code{Hayvan}'dır.

\begin{lstlisting}
#include <iostream>
#include <string>

// TABAN sinif
class Animal {
protected:                        // turetilmis siniflar erisebilir, disari erisemez
    std::string name_;
public:
    Animal(std::string name) : name_(std::move(name)) {}
    void breathe() const { std::cout << name_ << " nefes aliyor\n"; }
    const std::string& name() const { return name_; }
};

// TURETILMIS sinif: Hayvan'dan public kalitim
class Dog : public Animal {
    std::string breed_;
public:
    // Taban yapicisini uye baslatma listesinde CAGIR
    Dog(std::string name, std::string breed)
        : Animal(std::move(name)), breed_(std::move(breed)) {}

    void bark() const { std::cout << name_ << " (" << breed_ << "): Hav!\n"; }
};

int main() {
    Dog k("Karabas", "Kangal");
    k.breathe();     // TABAN'dan devralindi
    k.bark();        // KENDI metodu
    std::cout << k.name() << "\n";   // taban'in public metodu
    return 0;
}
\end{lstlisting}

\begin{center}
\begin{tabular}{@{}lp{9.2cm}@{}}
\toprule
\textbf{Kalıtım türü} & \textbf{Etkisi} \\
\midrule
\code{public} & Taban'ın public'i public, protected'ı protected kalır. (En yaygın, ``is-a''.) \\
\code{protected} & Taban'ın public ve protected üyeleri protected olur. \\
\code{private} & Taban'ın tüm üyeleri private olur. (``ile-gerçeklenmiştir''.) \\
\bottomrule
\end{tabular}
\end{center}

% =====================================================================
\gsec{Polimorfizm ve Sanal Fonksiyonlar}

Polimorfizm (``çok biçimlilik''), bir taban sınıf işaretçisi/referansı
üzerinden çağrılan bir fonksiyonun, \emph{çalışma zamanında} gerçek nesnenin
tipine göre doğru sürümünü çalıştırmasıdır. Bu, \code{virtual} anahtar
kelimesiyle sağlanır ve nesne yönelimli tasarımın en güçlü aracıdır.

\begin{lstlisting}
#include <iostream>
#include <vector>
#include <memory>

// Soyut taban sinif (abstract base class / arayuz)
class Shape {
public:
    // Saf sanal fonksiyon (= 0): govdesi yok, alt siniflar MECBUREN yazar
    virtual double area() const = 0;
    virtual void draw() const = 0;

    // SANAL YIKICI: taban isaretcisiyle silmede kritik!
    virtual ~Shape() = default;
};

class Circle : public Shape {
    double radius_;
public:
    Circle(double r) : radius_(r) {}
    double area() const override { return 3.14159 * radius_ * radius_; }
    void draw() const override { std::cout << "Daire (r=" << radius_ << ")\n"; }
};

class Rectangle : public Shape {
    double width_, height_;
public:
    Rectangle(double e, double b) : width_(e), height_(b) {}
    double area() const override { return width_ * height_; }
    void draw() const override { std::cout << "Dikdortgen " << width_ << "x" << height_ << "\n"; }
};

int main() {
    // Taban isaretcileriyle farkli turleri TEK kapta tut (polimorfizm)
    std::vector<std::unique_ptr<Shape>> shapes;
    shapes.push_back(std::make_unique<Circle>(2.0));
    shapes.push_back(std::make_unique<Rectangle>(3.0, 4.0));

    double total = 0;
    for (const auto& s : shapes) {
        s->draw();                     // dogru ciz() CALISMA ZAMANINDA secilir
        total += s->area();          // dogru alan() cagrilir
    }
    std::cout << "Toplam alan: " << total << "\n";
    return 0;
}   // unique_ptr'lar silinirken SANAL yikici sayesinde dogru yikici cagrilir
\end{lstlisting}

\uyari{\textbf{Sanal yıkıcı kritiktir:} Bir sınıf polimorfik olarak
kullanılacaksa (taban işaretçisiyle türetilmiş nesne tutulup silinecekse),
tabanın yıkıcısı \code{virtual} olmalıdır. Aksi halde \code{delete taban\_ptr}
yalnızca tabanın yıkıcısını çağırır, türetilmişinkini \emph{atlar} ---
kaynak sızıntısı ve tanımsız davranış. Kural: ``Sanal fonksiyonu olan her
sınıfın sanal yıkıcısı olmalıdır.''}

\bilgi{\code{override} anahtar kelimesi zorunlu değildir ama her zaman
kullanın: derleyici, gerçekten bir taban sanal fonksiyonunu geçersiz kılıp
kılmadığınızı kontrol eder. İmzada küçük bir hata (örneğin \code{const}
unutmak) yaparsanız, \code{override} olmadan sessizce \emph{yeni} bir
fonksiyon tanımlarsınız; \code{override} ile derleme hatası alır ve hatayı
anında görürsünüz. \code{final} ise bir sınıfın/fonksiyonun daha fazla
türetilmesini/geçersiz kılınmasını engeller.}

% #####################################################################
\gpart{3}{Bellek Modeli ve İşaretçiler}
% #####################################################################

% =====================================================================
\gsec{Programın Bellek Düzeni}

Bir C++ programı çalışırken belleği birkaç mantıksal bölgeye ayrılır. Bu
bölgeleri anlamak, kaynak yönetiminin temelidir; çünkü bir nesnenin
\emph{nerede} yaşadığı, ne kadar yaşadığını ve onu kimin temizlemesi
gerektiğini belirler.

\begin{center}
\begin{tabular}{@{}>{\raggedright\arraybackslash}p{3.0cm} >{\raggedright\arraybackslash}p{9.0cm}@{}}
\toprule
\textbf{Bölge} & \textbf{Açıklama} \\
\midrule
\textbf{Kod (text)} & Makine talimatları. Salt okunur, sabit boyut. \\
\textbf{Statik/Global} & Global ve \code{static} değişkenler. Program boyunca yaşar. \\
\textbf{Yığın (stack)} & Yerel değişkenler, fonksiyon çağrı çerçeveleri.
Otomatik ömür: kapsam bitince \emph{otomatik} yok edilir. Hızlı ama sınırlı. \\
\textbf{Öbek (heap)} & \code{new}/\code{malloc} ile ayrılan bellek. Ömrü
programcı yönetir. Esnek ama pahalı ve sızıntıya açık. \\
\bottomrule
\end{tabular}
\end{center}

\begin{lstlisting}
#include <iostream>

int global_value = 10;          // statik/global bolge

void func() {
    int local = 5;              // YIĞIN (stack): kapsam bitince otomatik silinir
    static int calls = 0;       // STATIK: cagrilar arasi degerini korur
    int* heap_val = new int(42); // ÖBEK (heap): ELLE silinmeli!

    ++calls;
    std::cout << "yerel=" << local << " cagri=" << calls
              << " obek=" << *heap_val << "\n";

    delete heap_val;             // AKSI HALDE BELLEK SIZINTISI
}   // 'yerel' burada otomatik yok olur; 'obekten' isaretcisi de,
    // ama isaret ettigi obek bellegi delete edilmezse sizardi.

int main() {
    func();
    func();
    return 0;
}
\end{lstlisting}

\bilgi{\textbf{Otomatik ömür (automatic storage)} C++'ın en güçlü özelliğidir:
yığındaki nesneler, kapsam (scope, yani \code{\{ \}}) bittiğinde derleyici
tarafından \emph{garantili} olarak yok edilir; siz hiçbir şey yapmazsınız.
Bu garanti, ileride göreceğimiz RAII deseninin temelidir. Öbekteki nesnelerde
ise böyle bir garanti yoktur; onları siz temizlemekle yükümlüsünüz.}

% =====================================================================
\gsec{new ve delete: Öbek Belleği Yönetimi}

\code{new} operatörü öbekte bellek ayırır \emph{ve} nesnenin yapıcısını
çağırır. \code{delete} ise yıkıcıyı çağırır \emph{ve} belleği geri verir.
Bu ikisi C'deki \code{malloc}/\code{free}'den farklıdır: onlar yapıcı/yıkıcı
çalıştırmaz, yalnızca ham bellek yönetir.

\begin{lstlisting}
#include <iostream>

struct Resource {
    int id;
    Resource(int i) : id(i) { std::cout << "Kaynak " << id << " olustu\n"; }
    ~Resource()             { std::cout << "Kaynak " << id << " yok oldu\n"; }
};

int main() {
    // Tekil nesne
    Resource* k = new Resource(1);   // bellek ayIr + yapIcI cagIr
    std::cout << "id = " << k->id << "\n";
    delete k;                    // yIkIcI cagIr + bellek geri ver

    // Dizi: new[] ve delete[] BIRLIKTE kullanIlmalI
    Resource* arr = new Resource[3]{ {10}, {20}, {30} };
    std::cout << "dizi[1].id = " << arr[1].id << "\n";
    delete[] arr;               // MUTLAKA delete[] -- 'delete' YANLIS olur

    return 0;
}
\end{lstlisting}

\uyari{\code{new[]} ile ayrılan bellek \code{delete[]} ile, \code{new} ile
ayrılan \code{delete} ile serbest bırakılmalıdır. \code{new[]} belleğini
düz \code{delete} ile silmek \textbf{tanımsız davranıştır} (undefined
behavior); genellikle yalnızca ilk elemanın yıkıcısı çağrılır ve bellek
bozulur. Modern C++'ta bu hatalardan kaçınmak için \emph{ham} \code{new}/\code{delete}
yerine akıllı işaretçileri ve konteynerleri (\code{std::vector}) tercih edin.}

\gsub{Bellek Sızıntıları ve İkili Silme}

En sık yapılan iki hata: belleği hiç silmemek (sızıntı) veya aynı belleği
iki kez silmek (double free).

\begin{lstlisting}
void leak() {
    int* p = new int(5);
    // ... delete unutuldu -> SIZINTI: bu bellek program bitene kadar kayIp
}

void double_delete() {
    int* p = new int(5);
    delete p;
    delete p;      // TANIMSIZ DAVRANIS: ayni bellegi iki kez silmek
    // Cozum: sildikten sonra p = nullptr; yapmak
    //   (nullptr'I delete etmek guvenlidir, hicbir sey yapmaz)
}

void dangling() {
    int* p = new int(5);
    delete p;
    std::cout << *p;   // SARKAN ISARETCI: silinen bellege erisim -> UB
    p = nullptr;       // dogru alIskanlIk
}
\end{lstlisting}

\ipucu{Bu üç sorunun (sızıntı, ikili silme, sarkan işaretçi) \emph{kalıcı}
çözümü, ham \code{new}/\code{delete}'i elle yönetmeyi bırakıp sahipliği
akıllı işaretçilere devretmektir. Kısım 2'de bunu göreceğiz. Kural:
``Modern C++ kodunda çıplak \code{new} ve \code{delete} neredeyse hiç
görünmemelidir.''}

% =====================================================================
\gsec{İşaretçiler ve Referanslar}

İşaretçi (pointer), başka bir nesnenin \emph{bellek adresini} tutan bir
değişkendir. Referans (reference) ise mevcut bir nesneye verilen
\emph{takma addır} (alias). İkisi de dolaylı erişim sağlar ama davranışları
önemli ölçüde farklıdır.

\begin{lstlisting}
#include <iostream>

int main() {
    int x = 10;

    // ISARETCI: x'in adresini tutar; yeniden atanabilir, null olabilir
    int* p = &x;       // p, x'i gosterir
    *p = 20;           // dereference: x'in degerini degistirir
    std::cout << "x = " << x << "\n";   // 20

    int y = 99;
    p = &y;            // isaretci baskasini gosterecek sekilde DEGISTIRILEBILIR
    std::cout << "*p = " << *p << "\n"; // 99

    // REFERANS: x'e takma ad; olusturuldugu anda baglanIr, degistirilemez
    int& r = x;        // r, x'in ta kendisidir
    r = 30;            // x'i degistirir (yeni degisken DEGIL)
    std::cout << "x = " << x << "\n";   // 30
    // int& r2;        // HATA: referans mutlaka baslangicta baglanmalI
    // r = y;          // bu y'yi r'ye atamaz; x = y anlamIna gelir!

    return 0;
}
\end{lstlisting}

\begin{center}
\begin{tabular}{@{}>{\raggedright\arraybackslash}p{4.4cm} >{\raggedright\arraybackslash}p{3.8cm} >{\raggedright\arraybackslash}p{3.8cm}@{}}
\toprule
\textbf{Özellik} & \textbf{İşaretçi (\code{T*})} & \textbf{Referans (\code{T\&})} \\
\midrule
Null olabilir mi? & Evet (\code{nullptr}) & Hayır \\
Yeniden bağlanır mı? & Evet, başkasını gösterebilir & Hayır, sabit bağ \\
Başlatma zorunlu mu? & Hayır (ama tehlikeli) & Evet \\
Aritmetik yapılır mı? & Evet (\code{p+1}) & Hayır \\
Adres alınır mı? & \code{\&p} işaretçinin adresi & \code{\&r} nesnenin adresi \\
Kullanım & Sahiplik, dizi, opsiyonel & Fonksiyon parametresi, takma ad \\
\bottomrule
\end{tabular}
\end{center}

\bilgi{\textbf{Ne zaman hangisi?} Bir fonksiyona ``değiştirilebilir ama
her zaman geçerli'' bir nesne geçiriyorsanız \textbf{referans} (\code{T\&})
kullanın. ``Opsiyonel'' (olmayabilir) veya ``yeniden yönlendirilebilir''
bir bağlantı gerekiyorsa \textbf{işaretçi} kullanın. Büyük nesneleri
kopyalamadan salt okunur geçirmek için \code{const T\&} idealdir.}

\gsub{const ve İşaretçiler}

\code{const} anahtar kelimesinin işaretçilerle konumu, neyin sabit
olduğunu değiştirir. Bu, C++'ta sık kafa karıştıran ama kritik bir konudur.
Okumak için işaretçiyi \emph{sağdan sola} okuyun.

\begin{lstlisting}
int value = 10, other = 20;

// 1) Gosterilen SABIT, isaretci degisebilir
const int* p1 = &value;     // "pointer to const int"
// *p1 = 5;                 // HATA: gosterilen degeri degistiremez
p1 = &other;                // OK: isaretci baskasini gosterebilir

// 2) Isaretci SABIT, gosterilen degisebilir
int* const p2 = &value;     // "const pointer to int"
*p2 = 5;                    // OK: gosterilen degeri degistirir
// p2 = &other;            // HATA: isaretci yeniden atanamaz

// 3) Her ikisi de SABIT
const int* const p3 = &value;
// *p3 = 5;                // HATA
// p3 = &baska;           // HATA

// Okuma kurali (sagdan sola):
//   const int* const p3
//   p3, const bir (isaretci) -> const int'e
\end{lstlisting}

\ipucu{Pratik kural: \code{const}'ı işaretçi yıldızının \emph{soluna}
yazarsanız \emph{gösterilen} sabittir (\code{const int* p}); \emph{sağına}
yazarsanız \emph{işaretçinin kendisi} sabittir (\code{int* const p}).
Salt okunur veri geçirmek için en sık kullanılan biçim birincisidir:
\code{const T*} veya \code{const T\&}.}

\gsub{Boş, Sarkan ve Yaban İşaretçiler}

\begin{lstlisting}
int* wild;                 // YABAN (wild): baslatilmadi, cop adres -> UB
int* null_ptr = nullptr;         // BOS (null): guvenli, "hicbir seyi gostermiyor"

if (null_ptr == nullptr)         // guvenli kontrol
    std::cout << "bos isaretci\n";

int* dangling = new int(5);
delete dangling;              // bellek geri verildi
// sarkan artik SARKAN (dangling): gecersiz bellegi gosteriyor
dangling = nullptr;           // dogru alIskanlIk: sildikten sonra sIfIrla

// nullptr her zaman C++11+ ile kullanIlmalI; eski 'NULL' veya '0' DEGIL
void f(int);
void f(char*);
// f(NULL);   // belirsiz olabilir (NULL genelde 0'dir -> f(int))
// f(nullptr); // net: f(char*) cagrIlIr
\end{lstlisting}

\uyari{\code{nullptr}, C++11 ile gelen tip güvenli boş işaretçi sabitidir.
Eski \code{NULL} makrosu genellikle tam sayı \code{0} olarak tanımlıdır ve
aşırı yüklenmiş fonksiyonlarda yanlış olanı seçebilir. Modern kodda her
zaman \code{nullptr} kullanın.}

% #####################################################################
\gpart{4}{RAII ve Akıllı İşaretçiler}
% #####################################################################

% =====================================================================
\gsec{RAII: Kaynak Edinimi Başlatmadır}

RAII (Resource Acquisition Is Initialization), C++'ın en önemli deyimidir.
Fikir basittir: bir \emph{kaynağı} (bellek, dosya, kilit, soket) bir
\emph{nesnenin} yapıcısında edin, yıkıcısında serbest bırak. Nesne yığında
yaşadığından, kapsam bittiğinde yıkıcı \emph{otomatik} ve \emph{garantili}
çalışır --- istisna atılsa bile. Böylece kaynak asla sızmaz.

\begin{lstlisting}
#include <iostream>

// RAII sarmalayIcI: bir dosya tanIticIsInI (handle) yonetir
class FileHolder {
    std::FILE* file_;
public:
    explicit FileHolder(const char* name, const char* mode)
        : file_(std::fopen(name, mode)) {
        if (!file_) throw std::runtime_error("Dosya acilamadi");
        std::cout << "Dosya acildi\n";
    }
    ~FileHolder() {                      // YIKICI: kaynagi serbest birak
        if (file_) { std::fclose(file_); std::cout << "Dosya kapandi\n"; }
    }
    std::FILE* get() const { return file_; }

    // Kopyalamayi engelle (kaynak tekil sahiplikte -- Kisim 3'te aciklanacak)
    FileHolder(const FileHolder&) = delete;
    FileHolder& operator=(const FileHolder&) = delete;
};

void write() {
    FileHolder f("/tmp/test.txt", "w");
    std::fputs("Merhaba RAII\n", f.get());
    // throw std::runtime_error("hata!");  // istisna olsa BILE dosya kapanIr
}   // kapsam bitti -> f'nin yIkIcIsI otomatik cagrIlIr -> dosya kapanIr

int main() {
    write();
    std::cout << "Bitti\n";
    return 0;
}
\end{lstlisting}

\bilgi{RAII'nin gücü \textbf{istisna güvenliğidir}: fonksiyonun ortasında
bir istisna atılsa bile, yığın çözülürken (stack unwinding) tüm yerel
nesnelerin yıkıcıları çalışır. Böylece dosya kapatma, kilit bırakma, bellek
serbest bırakma gibi temizlik işlemleri \emph{unutulamaz}. \code{std::vector},
\code{std::string}, \code{std::lock\_guard} ve akıllı işaretçiler --- hepsi
RAII uygular.}

% =====================================================================
\gsec{std::unique\_ptr: Tekil Sahiplik}

\code{std::unique\_ptr}, bir öbek nesnesinin \emph{tek} sahibidir.
Kopyalanamaz (iki sahip olamaz) ama \emph{taşınabilir} (sahiplik devri).
Kapsam bittiğinde işaret ettiği nesneyi otomatik siler. Ham işaretçiyle
aynı boyuttadır --- yani sıfır ek yük (zero overhead).

\begin{lstlisting}
#include <iostream>
#include <memory>

struct Object {
    int id;
    Object(int i) : id(i) { std::cout << "Nesne " << id << " olustu\n"; }
    ~Object()             { std::cout << "Nesne " << id << " silindi\n"; }
    void greet() const   { std::cout << "Ben " << id << "\n"; }
};

std::unique_ptr<Object> produce(int id) {
    // make_unique: istisna guvenli, tercih edilen olusturma yolu
    return std::make_unique<Object>(id);
}   // sahiplik cagIrana TASINIR (move) -- kopyalanmaz

int main() {
    std::unique_ptr<Object> a = std::make_unique<Object>(1);
    a->greet();                       // -> operatoru ham isaretci gibi calisir

    // KOPYALAMA YASAK -- derleme hatasi:
    // std::unique_ptr<Object> b = a;  // HATA: unique_ptr kopyalanamaz

    // TASIMA serbest -- sahiplik a'dan b'ye gecer
    std::unique_ptr<Object> b = std::move(a);
    if (!a) std::cout << "a artik bos (nullptr)\n";
    b->greet();

    // Fonksiyondan donen sahipligi al
    std::unique_ptr<Object> c = produce(2);

    // release / reset
    Object* raw = c.release();         // sahipligi BIRAK, ham isaretci al (silme SANA kalIr)
    delete raw;                       // artik elle silmelisin

    b.reset();                        // b'nin nesnesini simdi sil (Nesne 1 silinir)
    std::cout << "main sonu\n";
    return 0;
}   // program sonu: kalan tum unique_ptr'lar nesnelerini otomatik siler
\end{lstlisting}

\ipucu{Her zaman \code{std::make\_unique<T>(...)} kullanın, \code{unique\_ptr<T>(new T(...))}
değil. \code{make\_unique} daha kısadır, \code{T}'yi tekrar yazmaz ve
istisna güvenlidir. \code{unique\_ptr}, dizi için de çalışır:
\code{std::unique\_ptr<int[]> d = std::make\_unique<int[]>(100);} otomatik
olarak \code{delete[]} kullanır.}

\gsub{Özel Silici (Custom Deleter)}

\code{unique\_ptr}, belleği nasıl serbest bırakacağını özelleştirmenize
izin verir. Bu, \code{delete} dışında bir temizlik gerektiren kaynaklar
(C API tanıtıcıları, \code{fclose}, \code{free}) için idealdir.

\begin{lstlisting}
#include <cstdio>
#include <memory>

int main() {
    // FILE* icin fclose ile temizleyen unique_ptr
    auto deleter = [](std::FILE* f){ if (f) std::fclose(f); };
    std::unique_ptr<std::FILE, decltype(deleter)> file(
        std::fopen("/tmp/veri.txt", "w"), deleter);

    if (file) std::fputs("ozel silici ile\n", file.get());
    return 0;
}   // dosya kapsam disi -> lambda silici cagrIlIr -> fclose
\end{lstlisting}

% =====================================================================
\gsec{std::shared\_ptr ve std::weak\_ptr: Paylaşımlı Sahiplik}

\code{std::shared\_ptr}, bir nesnenin \emph{birden fazla} sahibi olmasına
izin verir. Bir \emph{referans sayacı} tutar: her kopya sayacı artırır, her
yok olma azaltır. Sayaç sıfıra düştüğünde nesne silinir. Bu esneklik bir
maliyetle gelir: sayaç için ek bir kontrol bloğu ve atomik güncellemeler.

\begin{lstlisting}
#include <iostream>
#include <memory>

struct Node {
    int value;
    Node(int v) : value(v) { std::cout << "Dugum " << v << " olustu\n"; }
    ~Node()                { std::cout << "Dugum " << value << " silindi\n"; }
};

int main() {
    std::shared_ptr<Node> a = std::make_shared<Node>(1);
    std::cout << "Sayac: " << a.use_count() << "\n";   // 1

    {
        std::shared_ptr<Node> b = a;    // KOPYA: sayac artar
        std::cout << "Sayac: " << a.use_count() << "\n";  // 2
        b->value = 100;                  // ayni nesne (a ile paylasImlI)
        std::cout << "a->deger = " << a->value << "\n";   // 100
    }   // b yok oldu -> sayac azalIr

    std::cout << "Sayac: " << a.use_count() << "\n";   // 1
    return 0;
}   // a yok oldu -> sayac 0 -> "Dugum 1 silindi"
\end{lstlisting}

\uyari{\code{shared\_ptr} her nesne için ayrı bir kontrol bloğu tahsis eder
ve sayaç güncellemeleri iş parçacığı güvenliği için atomiktir --- bu,
\code{unique\_ptr}'a göre gözle görülür bir ek yüktür. \textbf{Varsayılan
tercihiniz her zaman \code{unique\_ptr} olmalıdır}; yalnızca sahipliğin
gerçekten \emph{paylaşılması} gerektiğinde \code{shared\_ptr}'a geçin.
\code{make\_shared}, nesneyi ve kontrol bloğunu tek tahsiste birleştirerek
performansı iyileştirir.}

\gsub{Döngüsel Referans Sorunu ve weak\_ptr}

\code{shared\_ptr}'ın en büyük tuzağı \textbf{döngüsel referanstır}: iki
nesne birbirini \code{shared\_ptr} ile tutarsa, sayaçları asla sıfıra
düşmez ve ikisi de asla silinmez --- bir bellek sızıntısı. Çözüm,
döngüyü kıran taraflardan birinde \code{std::weak\_ptr} kullanmaktır.
\code{weak\_ptr}, nesneyi \emph{sahiplenmeden} gözlemler; sayacı etkilemez.

\begin{lstlisting}
#include <iostream>
#include <memory>

struct Child;
struct Parent {
    std::shared_ptr<Child> child;
    ~Parent() { std::cout << "Ebeveyn silindi\n"; }
};
struct Child {
    // DIKKAT: ebeveyni shared_ptr ile tutsaydIk DONGU olurdu.
    // weak_ptr ile sahiplenmeden gozlemliyoruz:
    std::weak_ptr<Parent> parent;
    ~Child() { std::cout << "Cocuk silindi\n"; }
};

int main() {
    auto e = std::make_shared<Parent>();
    auto c = std::make_shared<Child>();

    e->child = c;        // Ebeveyn -> Cocuk (shared, sahiplik)
    c->parent = e;      // Cocuk -> Ebeveyn (weak, gozlem) -> DONGU YOK

    // weak_ptr'I kullanmak icin once lock() ile shared_ptr'a cevir
    if (auto locked = c->parent.lock()) {
        std::cout << "Ebeveyn hala yasiyor, sayac: "
                  << locked.use_count() << "\n";
    }
    return 0;
}   // Her iki nesne de dogru sekilde silinir (dongu kIrIldI)
\end{lstlisting}

\bilgi{\code{weak\_ptr}, ``sahiplenmeyen gözlemci'' rolündedir. Nesneye
erişmek için önce \code{lock()} çağırıp geçici bir \code{shared\_ptr}
elde edersiniz; nesne çoktan silinmişse \code{lock()} boş bir
\code{shared\_ptr} döndürür ve siz güvenle kontrol edebilirsiniz. Ebeveyn--çocuk,
gözlemci (observer), önbellek gibi ilişkilerde döngüyü kırmak için idealdir.}

\begin{center}
\begin{tabular}{@{}>{\raggedright\arraybackslash}p{3.4cm} >{\raggedright\arraybackslash}p{8.6cm}@{}}
\toprule
\textbf{İşaretçi} & \textbf{Ne zaman kullanılır} \\
\midrule
\code{unique\_ptr} & \textbf{Varsayılan.} Tekil sahiplik. Sıfır ek yük. \\
\code{shared\_ptr} & Sahiplik gerçekten paylaşılmalıysa. Atomik sayaç, ek yük. \\
\code{weak\_ptr} & Sahiplenmeden gözlem; \code{shared\_ptr} döngülerini kırmak. \\
Ham \code{T*} & Sahiplenmeyen, salt gözlem/parametre. \emph{Asla \code{delete} etmeyin.} \\
\bottomrule
\end{tabular}
\end{center}

% #####################################################################
\gpart{5}{Özel Üye Fonksiyonlar --- Beşin Kuralı}
% #####################################################################

% =====================================================================
\gsec{Değer Kategorileri: lvalue ve rvalue}

Beşin Kuralını anlamak için önce \emph{değer kategorilerini} anlamak
gerekir. Basitçe: \textbf{lvalue}, bir adı/adresi olan, kalıcı bir nesnedir
(soluna atama yapılabilir). \textbf{rvalue}, geçici, isimsiz, birazdan yok
olacak bir değerdir (bir ifadenin sonucu, geçici nesne).

\begin{lstlisting}
#include <string>

std::string produce() { return "gecici"; }

int main() {
    std::string a = "merhaba";   // a bir LVALUE (isimli, kalIcI)
    std::string b = a;           // a lvalue -> KOPYA yapIlIr

    std::string c = produce();      // uret()'in dondurdugu gecici bir RVALUE
                                 // -> TASIMA (move) yapIlabilir (ucuz)
    std::string d = a + b;       // (a+b) gecici bir RVALUE

    // std::move: bir lvalue'yu rvalue'ymus GIBI davranmaya "iter"
    std::string e = std::move(a);  // a'nIn icerigi e'ye TASINIR
    // artik a "gecerli ama belirsiz" durumda -- yeniden atamadan kullanmayIn
    return 0;
}
\end{lstlisting}

\bilgi{\textbf{rvalue referansı} (\code{T\&\&}), yalnızca rvalue'lara
bağlanabilen özel bir referanstır ve C++11'in taşıma semantiğinin
temelidir. Bir nesne ``birazdan zaten yok olacaksa'', kaynaklarını
\emph{kopyalamak} yerine \emph{çalmak} (taşımak) çok daha ucuzdur. İşte
taşıma yapıcı ve taşıma atama operatörü tam olarak bunu yapar.}

\uyari{\code{std::move} aslında hiçbir şeyi \emph{taşımaz}! Yalnızca bir
statik dönüşümdür: lvalue'yu rvalue referansına çevirerek derleyiciye
``bu nesnenin kaynaklarını çalabilirsin'' der. Gerçek taşıma işini, çağrılan
taşıma yapıcı/atama operatörü yapar. \code{std::move}'dan sonra kaynak
nesne ``geçerli ama belirsiz'' durumdadır; yeniden atamadan içeriğine
güvenmeyin.}

% =====================================================================
\gsec{Yıkıcı, Kopya ve Taşıma --- Beş Fonksiyon Bir Arada}

Kaynak (öbek belleği, dosya, soket...) yöneten her sınıfın, nesnenin tüm
\emph{yaşam döngüsünü} kontrol eden \textbf{beş özel üye fonksiyonu} vardır.
Bunlar ayrı ayrı değil, bir \emph{aile} olarak düşünülmelidir; çünkü hepsi
aynı kaynağı doğru ele almak zorundadır. Derleyici bu beşini kendisi
üretebilir, ancak sınıf elle bir kaynağı (\code{new} ile ayrılmış bellek)
yönetiyorsa, üretilen varsayılanlar \textbf{yanlıştır}: işaretçiyi olduğu
gibi kopyalarlar (sığ kopya) ve aynı bellek iki kez silinir. Bu yüzden
beşini de bilinçli tanımlamak gerekir.

\begin{center}
\begin{tabular}{@{}c>{\raggedright\arraybackslash}p{3.5cm}>{\ttfamily\footnotesize}l>{\raggedright\arraybackslash}p{3.4cm}@{}}
\toprule
\# & \textbf{Fonksiyon} & \textrm{\textbf{İmza (T için)}} & \textbf{Ne zaman çağrılır} \\
\midrule
1 & Yıkıcı & \textasciitilde{}T() & Nesne yok olduğunda \\
2 & Kopya yapıcı & T(const T\&) & Yeni nesne, mevcuttan \emph{kopyayla} kurulur \\
3 & Kopya atama & T\& operator=(const T\&) & Var olan nesneye lvalue atanır \\
4 & Taşıma yapıcı & T(T\&\&) noexcept & Yeni nesne, rvalue'dan kurulur \\
5 & Taşıma atama & T\& operator=(T\&\&) noexcept & Var olan nesneye rvalue atanır \\
\bottomrule
\end{tabular}
\end{center}

\gsub{Beşi Bir Arada: Tam Çalışan Örnek}

Aşağıdaki \code{TextHolder} sınıfı, öbekte bir karakter dizisi (\code{char*})
yönetir ve beş özel üye fonksiyonun \emph{tümünü} tanımlar. Her fonksiyon
hangi işlemin tetiklendiğini ekrana yazar; böylece \code{main} içindeki her
satırın hangi fonksiyonu çağırdığını çıktıdan net görebilirsiniz.

\begin{lstlisting}
#include <iostream>
#include <cstring>
#include <utility>

class TextHolder {
    std::size_t size_;    // NOT: size_ once bildirilir ki data_'dan ONCE baslasin
    char* data_;          // sahip olunan obek kaynagi
public:
    // (0) Normal yapici -- kaynagi EDINIR (RAII)
    explicit TextHolder(const char* s = "")
        : size_(std::strlen(s)), data_(new char[size_ + 1]) {
        std::strcpy(data_, s);
        std::cout << "[yapici]         \"" << data_ << "\"\n";
    }

    // (1) YIKICI -- kaynagi SERBEST BIRAKIR
    ~TextHolder() {
        std::cout << "[yikici]         \"" << (data_ ? data_ : "(bos)") << "\"\n";
        delete[] data_;
    }

    // (2) KOPYA YAPICI -- yeni, BAGIMSIZ kaynak ayirip icerigi kopyalar (derin kopya)
    TextHolder(const TextHolder& other)
        : size_(other.size_), data_(new char[other.size_ + 1]) {
        std::strcpy(data_, other.data_);
        std::cout << "[kopya yapici]   \"" << data_ << "\"\n";
    }

    // (3) KOPYA ATAMA -- var olan nesneye derin kopya atar
    TextHolder& operator=(const TextHolder& other) {
        std::cout << "[kopya atama]    \"" << other.data_ << "\"\n";
        if (this != &other) {                          // kendine-atama korumasi
            char* fresh = new char[other.size_ + 1];   // 1) once yeniyi ayir
            std::strcpy(fresh, other.data_);           //    (istisna guvenligi)
            delete[] data_;                            // 2) sonra eskiyi birak
            data_ = fresh;                             // 3) yeni duruma gec
            size_ = other.size_;
        }
        return *this;                                  // zincirleme icin (a = b = c)
    }

    // (4) TASIMA YAPICI -- kaynagi CALAR, orijinali gecerli-bos birakir (noexcept)
    TextHolder(TextHolder&& other) noexcept
        : size_(other.size_), data_(other.data_) {     // isaretciyi DEVRAL
        other.data_ = nullptr;                         // orijinali BOSALT
        other.size_ = 0;
        std::cout << "[tasima yapici]  \"" << data_ << "\"\n";
    }

    // (5) TASIMA ATAMA -- eski kaynagi birak, yeniyi cal (noexcept)
    TextHolder& operator=(TextHolder&& other) noexcept {
        std::cout << "[tasima atama]   \"" << (other.data_ ? other.data_ : "") << "\"\n";
        if (this != &other) {
            delete[] data_;              // eski kaynagi birak
            size_ = other.size_;
            data_ = other.data_;         // yeni kaynagi CAL
            other.data_ = nullptr;       // orijinali bosalt
            other.size_ = 0;
        }
        return *this;
    }

    const char* get() const { return data_ ? data_ : "(bos)"; }
};

TextHolder makeTemp() { return TextHolder("gecici"); }   // gecici (rvalue) uretir

int main() {
    TextHolder a("A");            // (0) yapici
    TextHolder b = a;             // (2) kopya yapici  -> yeni nesne, a'dan
    TextHolder c("C");            // (0) yapici
    c = a;                        // (3) kopya atama   -> c zaten VAR
    TextHolder d = std::move(a);  // (4) tasima yapici -> a'nin kaynagi calinir
    TextHolder e("E");            // (0) yapici
    e = makeTemp();               // (5) tasima atama  -> sagda gecici (rvalue)

    std::cout << "d = " << d.get() << ", a = " << a.get() << "\n";
    return 0;
}   // kapsam sonu: e, d, c, b, a icin (1) yikici -- TERS sirada
\end{lstlisting}

\ornek[Program Çıktısı]{\ttfamily\small
[yapici]~~~~~~~~~"A"\\
{}[kopya yapici]~~~"A"\\
{}[yapici]~~~~~~~~~"C"\\
{}[kopya atama]~~~~"A"\\
{}[tasima yapici]~~"A"\\
{}[yapici]~~~~~~~~~"E"\\
{}[yapici]~~~~~~~~~"gecici"~~~~(makeTemp icindeki gecici)\\
{}[tasima atama]~~~"gecici"\\
{}[yikici]~~~~~~~~~"(bos)"~~~~~(gecici, tasindiktan sonra bos)\\
d = A, a = (bos)\\
{}[yikici]~~~~~~~~~"gecici"~~~(e)\\
{}[yikici]~~~~~~~~~"A"~~~~~~~~~(d)\\
{}[yikici]~~~~~~~~~"A"~~~~~~~~~(c)\\
{}[yikici]~~~~~~~~~"A"~~~~~~~~~(b)\\
{}[yikici]~~~~~~~~~"(bos)"~~~~~(a, kaynagi calinmisti)}

\gsub{1. Yıkıcı (Destructor)}

Yıkıcı, nesne yok olduğunda \emph{otomatik} çağrılır ve nesnenin sahip
olduğu kaynağı serbest bırakır. \code{\~T()} biçiminde yazılır; parametre
almaz, değer döndürmez, aşırı yüklenemez. Yukarıda \code{delete[] data\_}
ile öbek belleğini geri verir. Çıktıda görüldüğü gibi, yerel nesneler
\emph{kuruluşun tersi sırada} (e, d, c, b, a) yok edilir.

\uyari{Bir sınıf polimorfik kullanılacaksa (taban işaretçisiyle türetilmiş
nesne tutulup silinecekse), yıkıcısı \code{virtual} olmalıdır ---
\code{virtual \~Base() = default;}. Aksi halde \code{delete base\_ptr}
türetilmiş sınıfın yıkıcısını atlar ve kaynak sızar.}

\gsub{2. Kopya Yapıcı (Copy Constructor)}

Kopya yapıcı, \emph{yeni} bir nesneyi var olan başka bir nesnenin kopyası
olarak \emph{ilklendirir} (\code{TextHolder b = a;}). Kritik nokta:
\textbf{derin kopya} yapması gerekir --- yeni, bağımsız bir öbek ayırıp
içeriği oraya kopyalar. Böylece \code{a} ile \code{b} tamamen ayrı
belleklere sahip olur.

\uyari{\textbf{Sığ kopya tuzağı:} Kopya yapıcı yazmazsanız, derleyici
işaretçiyi olduğu gibi kopyalayan bir varsayılan üretir --- iki nesne
\emph{aynı} öbeği gösterir. Yıkıcıda önce biri, sonra diğeri \emph{aynı}
belleği siler: \textbf{ikili silme} (double free) ve çökme. Derin kopya
tam olarak bunu önler.}

\ornek[Sığ vs Derin Kopya]{\ttfamily\small
SIĞ (yanlış):\\
~~a.data\_ ---.\\
~~~~~~~~~~~~~~~\textgreater{} [ "A\textbackslash0" ]~~~~(TEK obek, iki sahip!)\\
~~b.data\_ ---'\\[4pt]
DERİN (doğru):\\
~~a.data\_ ---\textgreater{} [ "A\textbackslash0" ]\\
~~b.data\_ ---\textgreater{} [ "A\textbackslash0" ]~~~(iki AYRI obek)}

\gsub{3. Kopya Atama Operatörü (Copy Assignment)}

Kopya atama, \emph{zaten var olan} bir nesneye başka bir nesnenin değerini
\emph{atar} (\code{c = a;}). Kopya yapıcıdan farkı: hedef nesnenin
\emph{eski bir kaynağı vardır} ve önce onu temizlemek gerekir. Doğru bir
kopya atamada üç kritik nokta:

\begin{enumerate}
\item \textbf{Kendine atama koruması:} \code{if (this != \&other)} ---
\code{x = x} durumunda önce kendi kaynağını silip sonra ondan kopyalamaya
çalışmayı önler.
\item \textbf{İstisna güvenliği:} Önce yeni belleği ayırın, \emph{sonra}
eskiyi silin. Ayırma başarısız olursa (istisna), nesne eski geçerli
durumunda kalır.
\item \textbf{Zincirleme için \code{*this} döndürün:} \code{a = b = c}
ifadesinin çalışması için.
\end{enumerate}

\gsub{4. Taşıma Yapıcı (Move Constructor)}

Kopyalama pahalıdır: yeni öbek ayırıp içeriği baytı baytına kopyalar. Ama
kaynak nesne zaten \emph{geçici} (rvalue) ise ve birazdan yok olacaksa,
kaynağını \emph{kopyalamak} yerine \emph{çalmak} çok daha ucuzdur: sadece
işaretçiyi devralır, orijinali \code{nullptr} yaparsınız. Bu, C++11'in
taşıma semantiğidir ve modern C++ performansının temelidir.
\code{TextHolder d = std::move(a);} satırı bunu tetikler --- hiç bellek
kopyalanmaz, yalnızca işaretçi \code{a}'dan \code{d}'ye geçer.

\ipucu{Taşıma işlemlerini \textbf{her zaman \code{noexcept}} yapın.
\code{std::vector} yeniden boyutlanırken elemanlarını yalnızca taşıma
yapıcı \code{noexcept}'se \emph{taşır}; değilse güvenlik için
\emph{kopyalar}. Yani \code{noexcept} olmayan bir taşıma yapıcı, vector
büyüdüğünde sessizce performans kaybettirir.}

\gsub{5. Taşıma Atama Operatörü (Move Assignment)}

Taşıma atama, var olan bir nesneye bir rvalue atar (\code{e = makeTemp();}).
Önce kendi eski kaynağını bırakır, sonra sağdaki geçicinin kaynağını çalar
ve onu boşaltır. Kopya atamanın ucuz kardeşidir: bellek kopyalamak yerine
sahiplik devreder.

\bilgi{\textbf{Taşıma sonrası durum:} \code{std::move} bir şeyi
\emph{taşımaz}; yalnızca nesneyi rvalue'ya çevirerek ``kaynağını
çalabilirsin'' der. Gerçek taşımayı yukarıdaki fonksiyonlar yapar. Taşınan
kaynak nesne yok \emph{edilmez}; ``geçerli ama belirsiz'' durumda kalır.
Onu tutarlı bir boş duruma getirmek (\code{data\_ = nullptr}) şarttır ki
yıkıcısı \code{delete[] nullptr} (güvenli) çağırsın. Çıktıda \code{a}'nın
kaynağı \code{d}'ye taşındıktan sonra \code{a.get()} ``(bos)'' döndürür.}

\notk[Hangisi Ne Zaman?]{Bir ifadenin kopya mı taşıma mı tetikleyeceğini
sağ taraf belirler: sağ taraf bir \textbf{lvalue} (isimli, kalıcı nesne)
ise \emph{kopya}, bir \textbf{rvalue} (geçici, \code{std::move}'lu) ise
\emph{taşıma} çağrılır. \code{b = a} kopya, \code{b = std::move(a)} veya
\code{b = makeTemp()} taşımadır. Hedefin yeni mi (yapıcı) yoksa var olan mı
(atama) olduğu ise \emph{yapıcı} ile \emph{atama} operatörünü ayırır.}

% =====================================================================
\gsec{Üç, Beş ve Sıfır Kuralı}

\gsub{Üçün Kuralı (Rule of Three) --- C++98}

Eğer bir sınıfın şu üçünden \emph{birini} elle yazmanız gerekiyorsa,
büyük olasılıkla \emph{üçünü de} yazmanız gerekir: \textbf{yıkıcı, kopya
yapıcı, kopya atama operatörü}. Çünkü elle yönetilen bir kaynak (öbek
belleği vb.) varsa, üçü de o kaynağı doğru ele almalıdır.

\gsub{Beşin Kuralı (Rule of Five) --- C++11}

C++11 ile taşıma semantiği gelince, liste beşe çıktı. Elle kaynak yöneten
bir sınıf genellikle şu \textbf{beşini} tanımlamalıdır:

\begin{center}
\begin{tabular}{@{}c>{\raggedright\arraybackslash}p{5.2cm}>{\ttfamily}l@{}}
\toprule
\# & \textbf{Fonksiyon} & \textrm{\textbf{İmza}} \\
\midrule
1 & Yıkıcı & \textasciitilde{}T() \\
2 & Kopya yapıcı & T(const T\&) \\
3 & Kopya atama & T\& operator=(const T\&) \\
4 & Taşıma yapıcı & T(T\&\&) noexcept \\
5 & Taşıma atama & T\& operator=(T\&\&) noexcept \\
\bottomrule
\end{tabular}
\end{center}

\uyari{\textbf{Kritik kural:} Bu beş fonksiyondan herhangi birini elle
tanımlarsanız, derleyici diğer bazılarını otomatik üretmeyi \emph{durdurabilir}.
Örneğin bir yıkıcı yazarsanız, derleyici taşıma işlemlerini \emph{üretmez}
(ve kopya işlemleri ``kullanımdan kaldırılmış'' sayılır). Bu yüzden ya
hepsini bilinçli tanımlayın, ya da hiçbirini (aşağıdaki Sıfır Kuralı).}

\gsub{Sıfırın Kuralı (Rule of Zero) --- En İyi Yaklaşım}

Modern C++'ın \emph{tercih edilen} yaklaşımı, bu beş fonksiyonun
\textbf{hiçbirini} yazmamaktır! Kaynakları doğrudan yönetmek yerine,
zaten doğru yöneten türleri (\code{std::vector}, \code{std::string},
\code{std::unique\_ptr}) üye olarak kullanın. O zaman derleyicinin ürettiği
varsayılan özel üye fonksiyonlar \emph{zaten doğru} çalışır.

\begin{lstlisting}
#include <string>
#include <vector>
#include <memory>

// SIFIRIN KURALI: hicbir ozel uye fonksiyon yazmadIk!
// Ciplak new/delete yok; kaynak yoneten uyeler kendileri halleder.
class User {
    std::string name_;                    // kendi kopya/tasImasInI yonetir
    std::vector<int> scores_;            // kendi kopya/tasImasInI yonetir
    std::unique_ptr<int> secret_;          // tasInabilir, kopyalanamaz
public:
    User(std::string name)
        : name_(std::move(name)),
          secret_(std::make_unique<int>(42)) {}

    // Yikici, kopya/tasIma... HICBIRINI yazmadIk.
    // Derleyici hepsini DOGRU sekilde uretir:
    //  - kopya: unique_ptr kopyalanamadIgI icin bu sInIf da kopyalanamaz
    //  - tasIma: uyelerin tasImasInI cagIrIr (otomatik, dogru, noexcept)
};
\end{lstlisting}

\bilgi{\textbf{En iyi tavsiye:} Sınıflarınızı Sıfır Kuralına göre tasarlayın.
Ham kaynak yöneten (elle \code{new}/\code{delete} yapan) sınıflar yazmaktan
kaçının; bunun yerine kaynağı bir \code{unique\_ptr} veya konteyner içine
sarın. Bu sayede Beşin Kuralı'nın tüm karmaşıklığı ve hata riski ortadan
kalkar. Beşin Kuralı'nı yalnızca gerçekten düşük seviye bir kaynak
sarmalayıcı (örneğin bir C API tanıtıcısı) yazarken uygulayın.}

\gsub{= default ve = delete ile Açık Kontrol}

C++11, özel üye fonksiyonların üretimini açıkça kontrol etmenizi sağlar.
\code{= default} derleyiciden varsayılanı üretmesini ister; \code{= delete}
o fonksiyonu tamamen yasaklar.

\begin{lstlisting}
class Example {
public:
    Example() = default;                         // varsayIlan yapIcIyI uret

    // Kopyalamayi acikca YASAKLA (ornegin tekil sahiplik icin)
    Example(const Example&) = delete;
    Example& operator=(const Example&) = delete;

    // Tasimaya IZIN ver (varsayIlanI uret)
    Example(Example&&) noexcept = default;
    Example& operator=(Example&&) noexcept = default;

    virtual ~Example() = default;                // varsayIlan sanal yIkIcI
};
\end{lstlisting}

\ipucu{\code{= default}, ``derleyicinin ürettiği varsayılan davranış tam
istediğim şey, ama niyetimi açıkça belgeliyorum'' demektir --- ayrıca bazı
durumlarda (örneğin başka bir yapıcı tanımladığınızda kaybolan varsayılan
yapıcıyı geri getirmek için) gereklidir. \code{= delete} ise kopyalamayı,
istenmeyen dönüşümleri veya belirli aşırı yüklemeleri derleme zamanında
yasaklamanın temiz yoludur.}

% =====================================================================
\gsec{Copy-and-Swap Deyimi}

Copy-and-swap, kopya atama operatörünü hem \emph{istisna güvenli} hem
\emph{kendine-atama güvenli} hem de kısa yazmanın zarif bir yoludur. Fikir:
argümanı \emph{değer olarak} al (böylece kopya otomatik oluşur), sonra
mevcut nesneyle \code{swap} et. Bonus: tek bir atama operatörü hem kopya
hem taşıma atamasını halleder.

\begin{lstlisting}
#include <algorithm>
#include <cstring>
#include <utility>

class TextHolder {
    char* data_ = nullptr;
    std::size_t size_ = 0;
public:
    TextHolder(const char* s = "")
        : size_(std::strlen(s)), data_(new char[size_ + 1]) {
        std::strcpy(data_, s);
    }
    ~TextHolder() { delete[] data_; }

    TextHolder(const TextHolder& d)
        : size_(d.size_), data_(new char[d.size_ + 1]) {
        std::strcpy(data_, d.data_);
    }
    TextHolder(TextHolder&& d) noexcept { swap_with(d); }   // bos'tan basla, takas et

    // Uyeleri takas eden yardImcI (ADL icin ozgur swap kullanIr)
    void swap_with(TextHolder& d) noexcept {
        using std::swap;
        swap(data_, d.data_);
        swap(size_, d.size_);
    }

    // TEK atama operatoru: argumanI DEGER olarak alIr.
    // - lvalue verilirse: kopya yapIcI cagrIlIr
    // - rvalue verilirse: tasIma yapIcI cagrIlIr
    // Sonra sadece takas ederiz. Kendine-atama ve istisna guvenligi BEDAVA.
    TextHolder& operator=(TextHolder d) noexcept {
        swap_with(d);
        return *this;
    }   // 'd' (eski icerigimizi tasIyan) burada yok olur, eski kaynak temizlenir

    const char* get() const { return data_; }
};
\end{lstlisting}

\bilgi{Copy-and-swap'ın zarafeti: argüman \code{d} \emph{değer olarak}
alındığı için, çağıran lvalue geçtiyse kopya yapıcı, rvalue geçtiyse taşıma
yapıcı otomatik devreye girer. Böylece \emph{tek} bir \code{operator=} hem
kopya hem taşıma atamasını halleder. Kendine-atama kontrolüne gerek yoktur
(kopya zaten oluşmuştur) ve güçlü istisna güvenliği sağlanır (takas
\code{noexcept}'tir, kopya başarısız olursa nesne değişmeden kalır).}

% #####################################################################
\gpart{6}{Şablonlar ve Generic Programlama}
% #####################################################################

% =====================================================================
\gsec{Fonksiyon Şablonları}

Şablonlar (templates), C++'ın \emph{jenerik programlama} aracıdır: tek bir
kod parçasını \emph{farklı tipler} için yeniden yazmadan çalıştırmanızı
sağlar. Derleyici, kullanılan her tip için şablonu \emph{somutlaştırır}
(instantiate). Bu, \code{std::vector} veya \code{std::sort} gibi tüm STL'in
temelidir.

\begin{lstlisting}
#include <iostream>
#include <string>

// Fonksiyon sablonu: T herhangi bir tip olabilir
template <typename T>
T maximum(T a, T b) {
    return (a > b) ? a : b;
}

// Iki farkli tip parametresi
template <typename T, typename U>
auto add(T a, U b) -> decltype(a + b) {   // donus tipini cikart
    return a + b;
}

int main() {
    // Derleyici T'yi argumanlardan CIKARIR (template argument deduction)
    std::cout << maximum(3, 7) << "\n";              // T=int -> 7
    std::cout << maximum(2.5, 1.5) << "\n";          // T=double -> 2.5
    std::cout << maximum<std::string>("elma", "armut") << "\n";  // acik tip

    // Karisik tipler: int + double -> double
    std::cout << add(3, 4.5) << "\n";              // 7.5
    return 0;
}
\end{lstlisting}

\bilgi{Derleyici \code{maximum(3, 7)} çağrısını görünce \code{T = int} için
somut bir fonksiyon üretir; \code{maximum(2.5, 1.5)} için ayrı bir
\code{double} sürümü üretir. Bu ``sıfır maliyetli soyutlamadır'': elle
yazılmış tipe özel koddan hiçbir çalışma zamanı farkı yoktur. Şablonlar
tamamen derleme zamanında çözülür.}

% =====================================================================
\gsec{Sınıf Şablonları}

Sınıflar da şablonlaştırılabilir. Aşağıda basit ama tam çalışan, jenerik
bir yığın (stack) uygulaması var --- \code{std::stack}'in nasıl çalıştığını
gösterir.

\begin{lstlisting}
#include <iostream>
#include <vector>
#include <stdexcept>

// Herhangi bir tip T tutabilen jenerik yigin
template <typename T>
class Stack {
    std::vector<T> items_;
public:
    void push(const T& value) { items_.push_back(value); }

    void pop() {
        if (empty()) throw std::out_of_range("Bos yigin");
        items_.pop_back();
    }

    const T& top() const {
        if (empty()) throw std::out_of_range("Bos yigin");
        return items_.back();
    }

    bool empty() const { return items_.empty(); }
    std::size_t size() const { return items_.size(); }
};

int main() {
    Stack<int> numbers;
    numbers.push(10); numbers.push(20); numbers.push(30);
    std::cout << "Tepe: " << numbers.top() << ", boyut: " << numbers.size() << "\n";
    numbers.pop();
    std::cout << "Pop sonrasi tepe: " << numbers.top() << "\n";   // 20

    // Ayni sablon, farkli tiple
    Stack<std::string> words;
    words.push("merhaba"); words.push("dunya");
    std::cout << words.top() << "\n";     // dunya
    return 0;
}
\end{lstlisting}

% =====================================================================
\gsec{Değişken Sayıda Şablon (Variadic) ve Kat İfadeleri}

Değişken sayıda şablon parametresi (variadic templates), keyfi sayıda
argüman alan jenerik fonksiyonlar yazmanızı sağlar. C++17'nin \emph{kat
ifadeleri} (fold expressions) bunu son derece zarif kılar.

\begin{lstlisting}
#include <iostream>

// C++17 kat ifadesi: tum argumanlari topla
template <typename... Args>
auto sum(Args... args) {
    return (args + ...);        // kat ifadesi: ((a1 + a2) + a3) + ...
}

// Her argumani ekrana yazan variadic fonksiyon
template <typename... Args>
void print(const Args&... args) {
    ((std::cout << args << " "), ...);   // virgul kat ifadesi
    std::cout << "\n";
}

int main() {
    std::cout << sum(1, 2, 3, 4, 5) << "\n";          // 15
    std::cout << sum(1.5, 2.5) << "\n";               // 4.0
    print("Sayilar:", 1, 2.5, 'X', "bitti");          // karisik tipler
    return 0;
}
\end{lstlisting}

\ipucu{Kat ifadeleri (\code{(args + ...)}) C++17 ile geldi ve variadic
şablonları okunması zor özyinelemeli yardımcılardan kurtardı. \code{(args + ...)}
sağ kat, \code{(... + args)} sol kattır. Virgül operatörüyle
(\code{(ifade, ...)}) her argüman için bir yan etki (örneğin yazdırma)
çalıştırabilirsiniz. Bu, modern jenerik kodun temel araçlarından biridir.}

% #####################################################################
\gpart{7}{STL: Konteynerler, İteratörler ve Algoritmalar}
% #####################################################################

% =====================================================================
\gsec{Standart Konteynerler}

STL (Standard Template Library), C++'ın hazır veri yapıları ve
algoritmalarından oluşan çekirdeğidir. Konteynerler, verileri farklı
şekillerde saklar; her birinin kendine özgü performans karakteristiği vardır.

\begin{center}
\begin{tabular}{@{}>{\ttfamily}l>{\raggedright\arraybackslash}p{5.4cm}>{\raggedright\arraybackslash}p{3.6cm}@{}}
\toprule
\textrm{\textbf{Konteyner}} & \textbf{Yapı} & \textbf{Tipik Erişim} \\
\midrule
vector & Dinamik dizi (bitişik bellek) & Sonda \bigO{1}, indeks \bigO{1} \\
deque & Çift uçlu kuyruk & İki uçta \bigO{1} \\
list & Çift bağlı liste & Ekleme/silme \bigO{1} \\
map & Sıralı (kırmızı-siyah ağaç) & Arama \bigO{\log n} \\
unordered\_map & Hash tablosu & Arama ortalama \bigO{1} \\
set & Sıralı benzersiz küme & \bigO{\log n} \\
\bottomrule
\end{tabular}
\end{center}

\begin{lstlisting}
#include <iostream>
#include <vector>
#include <map>
#include <unordered_map>
#include <set>
#include <string>

int main() {
    // vector: en cok kullanilan konteyner (dinamik dizi)
    std::vector<int> numbers{5, 3, 8, 1};
    numbers.push_back(9);
    std::cout << "vector[2] = " << numbers[2] << ", boyut = " << numbers.size() << "\n";

    // map: sirali anahtar-deger (otomatik sirali gezilir)
    std::map<std::string, int> ages;
    ages["Ada"] = 36;
    ages["Alan"] = 41;
    ages["Grace"] = 45;
    for (const auto& [name, age] : ages)          // yapisal baglama (C++17)
        std::cout << name << ": " << age << "\n";  // alfabetik sirali

    // unordered_map: hash tabanli (siralamaz ama ortalama O(1))
    std::unordered_map<std::string, double> prices;
    prices["kalem"] = 5.0;
    prices["defter"] = 12.5;
    if (auto it = prices.find("kalem"); it != prices.end())   // if-init (C++17)
        std::cout << "kalem fiyati: " << it->second << "\n";

    // set: benzersiz, sirali
    std::set<int> unique{3, 1, 4, 1, 5, 9, 2, 6, 5};  // tekrarlar elenir
    std::cout << "Benzersiz eleman sayisi: " << unique.size() << "\n";
    return 0;
}
\end{lstlisting}

% =====================================================================
\gsec{İteratörler}

İteratör (iterator), bir konteynerin elemanları üzerinde gezinmek için
kullanılan, işaretçiyi genelleştiren bir soyutlamadır. \code{begin()} ilk
elemanı, \code{end()} son elemanın \emph{bir sonrasını} gösterir. Tüm STL
algoritmaları iteratörler üzerinden çalışır --- bu sayede algoritmalar
konteyner tipinden bağımsızdır.

\begin{lstlisting}
#include <iostream>
#include <vector>
#include <list>

int main() {
    std::vector<int> v{10, 20, 30, 40};

    // Klasik iterator dongusu
    for (std::vector<int>::iterator it = v.begin(); it != v.end(); ++it)
        std::cout << *it << " ";        // *it: iteratorun gosterdigi eleman
    std::cout << "\n";

    // auto ile cok daha kisa
    for (auto it = v.begin(); it != v.end(); ++it)
        std::cout << *it << " ";
    std::cout << "\n";

    // Ters iterator (sondan basa)
    for (auto it = v.rbegin(); it != v.rend(); ++it)
        std::cout << *it << " ";        // 40 30 20 10
    std::cout << "\n";

    // Iteratorler farkli konteynerlerde ayni sekilde calisir
    std::list<char> chars{'a', 'b', 'c'};
    for (auto it = chars.begin(); it != chars.end(); ++it)
        std::cout << *it << " ";
    std::cout << "\n";
    return 0;
}
\end{lstlisting}

\bilgi{\code{end()} \emph{geçerli bir elemanı göstermez}; son elemanın bir
sonrasını (``geçmiş-son'', past-the-end) gösteren bir nöbetçidir. Döngü
koşulu bu yüzden \code{it != end()} biçimindedir. Bir iteratörü
\code{end()}'e ulaştıktan sonra dereference etmek (\code{*it}) tanımsız
davranıştır. Ayrıca konteyneri değiştirmek (örneğin \code{vector}'e ekleme
yapmak) mevcut iteratörleri \emph{geçersiz kılabilir} --- buna dikkat edin.}

% =====================================================================
\gsec{Algoritmalar}

\code{<algorithm>} başlığı, iteratör aralıkları üzerinde çalışan onlarca
hazır algoritma sunar: sıralama, arama, dönüştürme, sayma, biriktirme.
Kendi döngülerinizi yazmak yerine bunları kullanmak hem daha kısa hem daha
az hatalıdır.

\begin{lstlisting}
#include <iostream>
#include <vector>
#include <algorithm>
#include <numeric>

int main() {
    std::vector<int> v{5, 2, 8, 1, 9, 3};

    // sort: sirala (varsayilan artan)
    std::sort(v.begin(), v.end());
    for (int x : v) std::cout << x << " ";      // 1 2 3 5 8 9
    std::cout << "\n";

    // Lambda ile ozel olcut: azalan sirala
    std::sort(v.begin(), v.end(), [](int a, int b){ return a > b; });
    for (int x : v) std::cout << x << " ";      // 9 8 5 3 2 1
    std::cout << "\n";

    // find: bir eleman ara
    if (auto it = std::find(v.begin(), v.end(), 5); it != v.end())
        std::cout << "5 bulundu, konum: " << (it - v.begin()) << "\n";

    // count_if: kosulu saglayanlari say
    int evenCount = std::count_if(v.begin(), v.end(),
                                  [](int x){ return x % 2 == 0; });
    std::cout << "Cift sayi adedi: " << evenCount << "\n";   // 2

    // transform: her elemani donustur (karesini al)
    std::vector<int> squares(v.size());
    std::transform(v.begin(), v.end(), squares.begin(),
                   [](int x){ return x * x; });

    // accumulate: hepsini birlestir (topla)
    int total = std::accumulate(v.begin(), v.end(), 0);
    std::cout << "Toplam: " << total << "\n";   // 28

    // max_element / min_element
    auto maxIt = std::max_element(v.begin(), v.end());
    std::cout << "En buyuk: " << *maxIt << "\n";  // 9
    return 0;
}
\end{lstlisting}

\ipucu{``Ham döngü yazmadan önce bir STL algoritması var mı?'' diye sorun.
\code{std::sort}, \code{std::find}, \code{std::accumulate}, \code{std::transform},
\code{std::count\_if}, \code{std::any\_of}/\code{all\_of} gibi algoritmalar
niyetinizi net ifade eder ve hata olasılığını azaltır. C++20 ile gelen
\emph{ranges} (\code{std::ranges::sort(v)}) bunları daha da kısaltır ---
artık \code{begin()/end()} çifti yazmanıza bile gerek kalmaz.}

% #####################################################################
\gpart{8}{Lambda ve Fonksiyonel Programlama}
% #####################################################################

% =====================================================================
\gsec{Lambda İfadeleri}

Lambda ifadeleri, kod içinde \emph{yerinde} isimsiz fonksiyonlar
tanımlamanızı sağlar. STL algoritmalarına ölçüt geçmek, geri çağırmalar
ve kısa yardımcı işlemler için idealdir. Genel biçim:
\code{[yakalama](parametreler) \{ gövde \}}.

\begin{lstlisting}
#include <iostream>
#include <vector>
#include <algorithm>

int main() {
    // En basit lambda
    auto greet = []{ std::cout << "Merhaba lambda!\n"; };
    greet();

    // Parametreli ve donusluu
    auto add = [](int a, int b) { return a + b; };
    std::cout << add(3, 4) << "\n";     // 7

    // YAKALAMA (capture): disaridaki degiskenleri kullan
    int factor = 10;
    auto multiply = [factor](int x) { return x * factor; };  // deger ile yakala
    std::cout << multiply(5) << "\n";   // 50

    // Referansla yakalama: disaridaki degiskeni DEGISTIR
    int counter = 0;
    auto increment = [&counter]{ ++counter; };
    increment(); increment(); increment();
    std::cout << "Sayac: " << counter << "\n";   // 3

    // [=] tumunu degerle, [&] tumunu referansla yakalar
    int a = 1, b = 2;
    auto sumAll = [=]{ return a + b; };      // a,b kopyalanir
    std::cout << sumAll() << "\n";           // 3

    // mutable: deger-yakalanan kopyayi lambda icinde degistirmeye izin ver
    auto gen = [n = 0]() mutable { return ++n; };   // baslatmali yakalama (C++14)
    std::cout << gen() << gen() << gen() << "\n";    // 123

    // Jenerik lambda (C++14): auto parametre
    auto genericMax = [](auto x, auto y){ return x > y ? x : y; };
    std::cout << genericMax(3, 7) << " " << genericMax(2.5, 1.5) << "\n";
    return 0;
}
\end{lstlisting}

\uyari{Referansla yakalamada (\code{[\&]}) ömür yönetimine dikkat edin:
lambda, yakaladığı referanstan \emph{daha uzun} yaşarsa (örneğin bir
fonksiyondan döndürülüp sonra çağrılırsa), sarkan referans oluşur ve
tanımsız davranışa yol açar. Bir lambda'yı içinde tanımlandığı kapsamdan
dışarı taşıyacaksanız, referans yerine \emph{değerle} (\code{[=]} veya
açık kopya) yakalayın.}

% =====================================================================
\gsec{std::function ve Fonksiyon Nesneleri}

\code{std::function}, herhangi bir çağrılabilir nesneyi (lambda, serbest
fonksiyon, functor, üye fonksiyon) tek tip bir sarmalayıcıda tutan
genel amaçlı bir tiptir. Çağrılabilir nesneleri veri gibi saklamak,
geçirmek ve kaydetmek (örneğin geri çağırma tabloları) için kullanılır.

\begin{lstlisting}
#include <iostream>
#include <functional>
#include <vector>
#include <map>

int freeFunction(int x) { return x + 1; }

int main() {
    // std::function farkli cagrilabilirleri AYNI tipte tutar
    std::function<int(int)> f;

    f = freeFunction;                        // serbest fonksiyon
    std::cout << f(10) << "\n";              // 11

    f = [](int x){ return x * x; };          // lambda
    std::cout << f(5) << "\n";               // 25

    // Cagrilabilirleri bir kapta sakla
    std::vector<std::function<int(int)>> operations = {
        [](int x){ return x + 1; },
        [](int x){ return x * 2; },
        [](int x){ return x * x; }
    };
    for (auto& op : operations)
        std::cout << op(5) << " ";           // 6 10 25
    std::cout << "\n";

    // Komut tablosu: string -> islev (dispatch table deseni)
    std::map<std::string, std::function<double(double,double)>> calc = {
        {"topla",   [](double a, double b){ return a + b; }},
        {"cikar",   [](double a, double b){ return a - b; }},
        {"carp",    [](double a, double b){ return a * b; }}
    };
    std::cout << "3 carp 4 = " << calc["carp"](3, 4) << "\n";   // 12
    return 0;
}
\end{lstlisting}

\ipucu{\code{std::function} esnektir ama küçük bir çalışma zamanı maliyeti
(tip silme, olası yığın tahsisi, dolaylı çağrı) taşır. Performans kritik
sıcak döngülerde, tipi bilinen bir lambda'yı doğrudan (veya bir şablon
parametresi olarak) kullanmak daha hızlıdır. \code{std::function}'ı asıl
gücü olan \emph{heterojen çağrılabilirleri saklama} senaryolarında
(geri çağırma listeleri, komut tabloları) kullanın.}

% #####################################################################
\gpart{9}{İstisna Yönetimi (Exceptions)}
% #####################################################################

% =====================================================================
\gsec{try, catch ve throw}

İstisnalar (exceptions), çalışma zamanı hatalarını --- fonksiyonun normal
dönüş değerini kirletmeden --- bildirmenin ve ele almanın yapılandırılmış
bir yoludur. Bir hata oluştuğunda \code{throw} ile istisna \emph{fırlatılır};
çağrı yığını, onu yakalayacak bir \code{catch} bulunana kadar \emph{çözülür}
(stack unwinding) --- bu sırada tüm yerel nesnelerin yıkıcıları çalışır.

\begin{lstlisting}
#include <iostream>
#include <stdexcept>
#include <string>

double divide(double a, double b) {
    if (b == 0)
        throw std::invalid_argument("Sifira bolme!");   // istisna firlat
    return a / b;
}

int main() {
    try {
        std::cout << divide(10, 2) << "\n";   // 5
        std::cout << divide(10, 0) << "\n";   // istisna firlatir
        std::cout << "Bu satir CALISMAZ\n";
    }
    catch (const std::invalid_argument& e) {  // ozel istisna tipini yakala
        std::cout << "Gecersiz arguman: " << e.what() << "\n";
    }
    catch (const std::exception& e) {          // genel taban: her std istisnasi
        std::cout << "Genel hata: " << e.what() << "\n";
    }
    catch (...) {                              // her seyi yakalayan son savunma
        std::cout << "Bilinmeyen hata\n";
    }

    std::cout << "Program devam ediyor\n";     // catch sonrasi normal akis
    return 0;
}
\end{lstlisting}

\bilgi{\code{catch} blokları \emph{en özelden en genele} sıralanmalıdır:
önce \code{std::invalid\_argument}, sonra taban sınıfı \code{std::exception},
en sonda \code{catch(...)}. İstisnaları \textbf{her zaman referansla}
(\code{const std::exception\&}) yakalayın; değerle yakalamak nesne
kırpılmasına (slicing) yol açar ve polimorfik \code{what()} mesajını
kaybettirir. \code{<stdexcept>} başlığı hazır istisna tipleri sunar:
\code{logic\_error}, \code{runtime\_error}, \code{out\_of\_range},
\code{invalid\_argument} vb.}

% =====================================================================
\gsec{Özel İstisnalar ve İstisna Güvenliği}

Kendi istisna tiplerinizi \code{std::exception}'dan türeterek tanımlayın.
Ayrıca, istisnalar RAII ile birlikte kullanıldığında (Kısım 4) kod
\emph{istisna güvenli} olur: bir istisna atıldığında kaynaklar otomatik
temizlenir.

\begin{lstlisting}
#include <iostream>
#include <stdexcept>
#include <string>
#include <memory>

// Ozel istisna tipi: std::runtime_error'dan turet
class NetworkError : public std::runtime_error {
    int errorCode_;
public:
    NetworkError(const std::string& message, int code)
        : std::runtime_error(message), errorCode_(code) {}
    int code() const { return errorCode_; }
};

void connect(bool fail) {
    // RAII kaynagi: istisna atilsa BILE otomatik temizlenir
    auto resource = std::make_unique<int>(42);

    if (fail)
        throw NetworkError("Baglanti reddedildi", 503);

    std::cout << "Baglanti basarili, kaynak = " << *resource << "\n";
}   // 'resource' kapsam disinda otomatik silinir (istisna olsa da olmasa da)

int main() {
    try {
        connect(false);   // basarili
        connect(true);    // istisna firlatir
    }
    catch (const NetworkError& e) {
        std::cout << "Ag hatasi [" << e.code() << "]: " << e.what() << "\n";
    }
    return 0;
}
\end{lstlisting}

\uyari{\textbf{Yıkıcılardan istisna atmayın.} C++11'den beri yıkıcılar
örtük olarak \code{noexcept}'tir; bir yıkıcı istisna atarsa (özellikle
yığın çözülürken başka bir istisna zaten aktifse) program \code{std::terminate}
ile derhal sonlanır. Ayrıca istisnaları normal akış kontrolü için
kullanmayın; onlar \emph{istisnai} (beklenmedik hata) durumlar içindir.
Sık gerçekleşen ``hata'' durumları için \code{std::optional} veya C++23
\code{std::expected} daha uygundur.}

% #####################################################################
\gpart{10}{Önemli Anahtar Kelimeler}
% #####################################################################

% =====================================================================
\gsec{explicit: İstenmeyen Dönüşümleri Engelleme}

Tek argümanlı bir yapıcı, C++'ta \emph{örtük (implicit) dönüşüm} olarak
kullanılabilir --- bu genellikle istenmeyen ve şaşırtıcı sonuçlar doğurur.
\code{explicit} anahtar kelimesi bu örtük dönüşümü yasaklar.

\begin{lstlisting}
#include <iostream>

class Meter {
    double value_;
public:
    // explicit OLMASAYDI: 'Metre m = 5.0;' ve hatta fonksiyona '5.0'
    // gecmek sessizce Metre'ye donusurdu -- tehlikeli.
    explicit Meter(double d) : value_(d) {}
    double value() const { return value_; }
};

void print(Meter m) { std::cout << m.value() << " metre\n"; }

int main() {
    Meter a(5.0);            // OK: dogrudan baslatma
    // Meter b = 5.0;        // HATA (explicit): ortuk donusum yasak
    Meter c = Meter(5.0);    // OK: acik donusum
    // print(10.0);         // HATA: 10.0 sessizce Metre'ye donmez
    print(Meter(10.0));     // OK: acik
    return 0;
}
\end{lstlisting}

\ipucu{\textbf{Kural:} Tek argümanla çağrılabilen yapıcıları --- özellikle
bir birim/tip sarmalayıcısıysa --- neredeyse her zaman \code{explicit} yapın.
Örtük dönüşümü yalnızca gerçekten istediğinizde (örneğin \code{std::string}'in
\code{const char*}'tan dönüşümü) açık bırakın. C++'ta ``varsayılan olarak
\code{explicit}'' iyi bir alışkanlıktır.}

% =====================================================================
\gsec{const ve constexpr}

\code{const}, ``bu değer değiştirilemez'' der (çalışma zamanı sabiti).
\code{constexpr}, ``bu değer/fonksiyon derleme zamanında hesaplanabilir''
der. \code{constexpr} çok daha güçlüdür ve performans için değerlidir.

\begin{lstlisting}
#include <iostream>
#include <array>

// constexpr fonksiyon: derleme zamanInda hesaplanabilir
constexpr int factorial(int n) {
    return (n <= 1) ? 1 : n * factorial(n - 1);
}

class Circle {
    double radius_;
public:
    constexpr Circle(double r) : radius_(r) {}

    // const uye fonksiyon: nesneyi DEGISTIRMEZ (this uzerinden yazamaz)
    double area() const { return 3.14159 * radius_ * radius_; }
};

int main() {
    const int x = 10;              // calisma zamanI sabiti
    // x = 20;                     // HATA: const

    constexpr int f5 = factorial(5);   // DERLEME zamanInda hesaplanIr: 120
    std::array<int, factorial(4)> arr; // constexpr, sablon argumanI olabilir (24)
    std::cout << "5! = " << f5 << ", dizi boyutu = " << arr.size() << "\n";

    const Circle d(2.0);
    std::cout << "Alan = " << d.area() << "\n";   // const nesnede const uye cagrIlIr
    return 0;
}
\end{lstlisting}

\bilgi{\textbf{const üye fonksiyon} (\code{alan() const}), o fonksiyonun
nesnenin durumunu değiştirmeyeceğine söz verir; bu sayede \code{const}
nesneler üzerinde çağrılabilir. Nesnenin durumunu değiştirmeyen \emph{her}
üye fonksiyonu \code{const} yapmak (``const doğruluğu'', const-correctness)
iyi bir tasarım disiplinidir. C++20'de \code{consteval} (mutlaka derleme
zamanı) ve \code{constinit} de eklendi.}

% =====================================================================
\gsec{noexcept: İstisna Atmama Garantisi}

\code{noexcept}, bir fonksiyonun istisna \emph{atmayacağını} bildirir. Bu
hem bir optimizasyon ipucudur hem de --- daha önce gördüğümüz gibi ---
taşıma semantiği için kritik bir sözleşmedir. \code{noexcept} bir fonksiyon
yine de istisna atarsa, program derhal \code{std::terminate} ile sonlanır.

\begin{lstlisting}
#include <iostream>
#include <vector>

int safe_add(int a, int b) noexcept {   // asla istisna atmaz
    return a + b;
}

// Tasima islemleri noexcept OLMALI (vector optimizasyonu icin)
struct Movable {
    int* data;
    Movable(Movable&& d) noexcept : data(d.data) { d.data = nullptr; }
    Movable& operator=(Movable&&) noexcept = default;
    ~Movable() { delete data; }
    Movable() : data(new int(0)) {}
};

int main() {
    // noexcept operatoru: bir ifadenin noexcept olup olmadIgInI sorgular
    std::cout << std::boolalpha
              << noexcept(safe_add(1, 2)) << "\n";   // true

    // vector, noexcept tasIma yapIcIsI olan tipleri buyurken TASIR (kopyalamaz)
    std::vector<Movable> v;
    v.reserve(2);
    v.emplace_back();
    v.emplace_back();   // yeniden tahsiste noexcept tasIma kullanIlIr
    return 0;
}
\end{lstlisting}

\uyari{Bir fonksiyonu \code{noexcept} işaretleyip sonra istisna atarsanız,
\code{catch} ile yakalayamazsınız --- program \code{std::terminate}
çağırarak anında ölür. Bu yüzden \code{noexcept}'i yalnızca gerçekten
istisna atmayacağından \emph{emin} olduğunuz fonksiyonlarda kullanın.
Yıkıcılar C++11'den beri örtük olarak \code{noexcept}'tir; yıkıcıdan
istisna atmak neredeyse her zaman bir tasarım hatasıdır.}

% =====================================================================
\gsec{static, mutable ve friend}

\gsub{static Üyeler}

\begin{lstlisting}
#include <iostream>

class Counter {
    static int total_;      // TUM nesneler icin ORTAK tek bir degisken
    int id_;
public:
    Counter() : id_(++total_) {}
    static int count() { return total_; }   // static uye fonksiyon (this yok)
    int id() const { return id_; }
};
int Counter::total_ = 0;      // static uye sInIf dISInda TANIMLANMALI

int main() {
    Counter a, b, c;
    std::cout << "Toplam nesne: " << Counter::count() << "\n";  // 3
    std::cout << "c'nin id'si: " << c.id() << "\n";             // 3
    return 0;
}
\end{lstlisting}

\gsub{mutable Üyeler}

\code{mutable}, bir üyenin \code{const} bir nesnede bile değiştirilebilmesini
sağlar. Genellikle önbellekleme, kilit veya sayaç gibi ``mantıksal durumu''
değiştirmeyen yardımcı üyeler için kullanılır.

\begin{lstlisting}
#include <iostream>

class Calculator {
    mutable int callCount_ = 0;   // const uyede bile degisebilir
    int value_;
public:
    Calculator(int d) : value_(d) {}

    // const fonksiyon ama mutable uyeyi degistirebilir
    int square() const {
        ++callCount_;             // mutable sayesinde OK
        return value_ * value_;
    }
    int callCount() const { return callCount_; }
};

int main() {
    const Calculator h(5);         // const NESNE
    std::cout << h.square() << " " << h.square() << "\n";
    std::cout << "Cagri: " << h.callCount() << "\n";   // 2
    return 0;
}
\end{lstlisting}

\gsub{friend Fonksiyonlar}

\code{friend}, bir fonksiyona veya sınıfa, başka bir sınıfın \code{private}
üyelerine erişim ayrıcalığı verir. En yaygın kullanımı, birazdan göreceğimiz
akış operatörünü (\code{operator<<}) aşırı yüklemektir.

\begin{lstlisting}
#include <iostream>

class Point {
    int x_, y_;      // private
public:
    Point(int x, int y) : x_(x), y_(y) {}
    // friend: bu serbest fonksiyon private uyelere erisebilir
    friend std::ostream& operator<<(std::ostream& os, const Point& p);
};

std::ostream& operator<<(std::ostream& os, const Point& p) {
    return os << "(" << p.x_ << ", " << p.y_ << ")";   // private erisim
}

int main() {
    Point p(3, 4);
    std::cout << p << "\n";     // (3, 4)
    return 0;
}
\end{lstlisting}

% #####################################################################
\gpart{11}{Operatör Aşırı Yükleme}
% #####################################################################

% =====================================================================
\gsec{Operatör Aşırı Yüklemenin Temelleri}

C++, kendi tiplerinizin \code{+}, \code{==}, \code{[]}, \code{<<} gibi
operatörlerle nasıl davranacağını tanımlamanıza izin verir. Bu, kendi
tiplerinizi yerleşik tipler kadar doğal kullanmanızı sağlar. Ancak
\emph{dikkatli} kullanılmalıdır: operatörler sezgisel anlamlarını korumalıdır
(\code{+} toplama gibi davranmalı, sürpriz yapmamalı).

Bir matematiksel vektör (2B) sınıfı üzerinden operatörlerin çoğunu görelim:

\begin{lstlisting}
#include <iostream>

class Vec2 {
    double x_, y_;
public:
    Vec2(double x = 0, double y = 0) : x_(x), y_(y) {}
    double x() const { return x_; }
    double y() const { return y_; }

    // --- Bilesik atama (uye fonksiyon olarak yazmak GELENEKTIR) ---
    Vec2& operator+=(const Vec2& r) { x_ += r.x_; y_ += r.y_; return *this; }
    Vec2& operator-=(const Vec2& r) { x_ -= r.x_; y_ -= r.y_; return *this; }
    Vec2& operator*=(double s)      { x_ *= s;    y_ *= s;    return *this; }

    // --- Birli (unary) operatorler ---
    Vec2 operator-() const { return Vec2(-x_, -y_); }   // negatif

    // --- Karsilastirma (uye) ---
    bool operator==(const Vec2& r) const { return x_ == r.x_ && y_ == r.y_; }
    bool operator!=(const Vec2& r) const { return !(*this == r); }

    friend std::ostream& operator<<(std::ostream&, const Vec2&);
};

// --- Ikili aritmetik: SERBEST fonksiyon olarak yazmak GELENEKTIR ---
// (Simetri icin: 'a + b' ve donusumler her iki operand icin de calIsIr)
Vec2 operator+(Vec2 l, const Vec2& r) { return l += r; }   // l KOPYA -> += -> don
Vec2 operator-(Vec2 l, const Vec2& r) { return l -= r; }
Vec2 operator*(Vec2 v, double s)      { return v *= s; }
Vec2 operator*(double s, Vec2 v)      { return v *= s; }   // s * v de calIssIn

std::ostream& operator<<(std::ostream& os, const Vec2& v) {
    return os << "(" << v.x_ << ", " << v.y_ << ")";
}

int main() {
    Vec2 a(1, 2), b(3, 4);
    std::cout << "a + b   = " << (a + b) << "\n";     // (4, 6)
    std::cout << "b - a   = " << (b - a) << "\n";     // (2, 2)
    std::cout << "a * 2   = " << (a * 2) << "\n";     // (2, 4)
    std::cout << "3 * a   = " << (3 * a) << "\n";     // (3, 6)
    std::cout << "-a      = " << (-a) << "\n";        // (-1, -2)
    std::cout << "a == b? " << std::boolalpha << (a == b) << "\n"; // false

    a += b;                                           // bilesik atama
    std::cout << "a += b  -> " << a << "\n";          // (4, 6)
    return 0;
}
\end{lstlisting}

\bilgi{\textbf{Üye mi, serbest fonksiyon mu?} Genel kurallar: bileşik atama
(\code{+=}, \code{-=}) ve nesnenin durumunu değiştiren operatörler
\textbf{üye} olmalıdır. Simetrik ikili operatörler (\code{+}, \code{-}, \code{==})
\textbf{serbest fonksiyon} olmalıdır --- böylece sol operand da örtük
dönüşüme uğrayabilir. Kalıplaşmış deyim: \code{operator+}'yı, üye
\code{operator+=} cinsinden yazın (yukarıdaki gibi). Bu, kod tekrarını
önler ve tutarlılığı garanti eder.}

% =====================================================================
\gsec{Karşılaştırma ve Uzay Gemisi Operatörü (C++20)}

C++20 öncesinde, bir tipi tam sıralanabilir yapmak için altı operatörü
(\code{==}, \code{!=}, \code{<}, \code{<=}, \code{>}, \code{>=}) elle
yazmanız gerekiyordu. C++20'nin \textbf{üç yönlü karşılaştırma operatörü}
(\code{<=>}, ``spaceship'') bunların hepsini tek satırda üretir.

\begin{lstlisting}
#include <compare>
#include <iostream>

struct Version {
    int major, minor, patch;

    // Tek satIr: derleyici ==, !=, <, <=, >, >= HEPSINI uretir
    // Uyeler sIrayla (once ana, sonra ara, sonra yama) karsIlastIrIlIr
    auto operator<=>(const Version&) const = default;
};

int main() {
    Version a{1, 2, 0}, b{1, 3, 0};

    std::cout << std::boolalpha;
    std::cout << "a < b : " << (a < b) << "\n";    // true (ara: 2 < 3)
    std::cout << "a == b: " << (a == b) << "\n";   // false
    std::cout << "a >= b: " << (a >= b) << "\n";   // false

    // <=> dogrudan kullanImI: sIralama sonucu doner
    auto result = a <=> b;
    if (result < 0)      std::cout << "a kucuk\n";
    else if (result > 0) std::cout << "a buyuk\n";
    else                std::cout << "esit\n";
    return 0;
}
\end{lstlisting}

\begin{lstlisting}[style=shstyle]
g++ -std=c++20 spaceship.cpp -o spaceship
\end{lstlisting}

\ipucu{\code{auto operator<=>(const T\&) const = default;} tek satırı, tipi
tam sıralanabilir yapar ve \code{std::sort}, \code{std::set}, \code{std::map}
ile doğrudan kullanılabilir hale getirir. Dönüş tipi karşılaştırma kategorisini
belirtir: \code{std::strong\_ordering} (tam sıra), \code{std::weak\_ordering}
(eşdeğerlik), \code{std::partial\_ordering} (örn. \code{NaN} içeren float).
\code{= default} ile derleyici en uygununu seçer.}

% =====================================================================
\gsec{Özel Operatörler: [], (), ve Dönüşüm}

\gsub{Alt Simge Operatörü []}

\begin{lstlisting}
#include <iostream>
#include <stdexcept>

class Array {
    int* data_;
    std::size_t size_;
public:
    Array(std::size_t n) : data_(new int[n]{}), size_(n) {}
    ~Array() { delete[] data_; }
    Array(const Array&) = delete;             // basitlik icin kopyayI kapat
    Array& operator=(const Array&) = delete;

    // Iki asIrI yukleme: degistirilebilir ve salt-okunur (const) erisim
    int& operator[](std::size_t i) {
        if (i >= size_) throw std::out_of_range("indeks");
        return data_[i];                    // referans dondurur -> a[i] = 5 calIsIr
    }
    const int& operator[](std::size_t i) const {
        if (i >= size_) throw std::out_of_range("indeks");
        return data_[i];
    }
    std::size_t size() const { return size_; }
};

int main() {
    Array a(5);
    for (std::size_t i = 0; i < a.size(); ++i)
        a[i] = static_cast<int>(i * i);     // operator[] (non-const) -> atama
    for (std::size_t i = 0; i < a.size(); ++i)
        std::cout << a[i] << " ";           // 0 1 4 9 16
    std::cout << "\n";
    return 0;
}
\end{lstlisting}

\gsub{Fonksiyon Çağrı Operatörü () --- Fonksiyon Nesneleri}

\code{operator()} aşırı yüklendiğinde, nesne bir fonksiyon gibi
\emph{çağrılabilir} (functor / function object) hale gelir. STL algoritmaları
ve lambdaların temelinde bu yatar.

\begin{lstlisting}
#include <iostream>
#include <vector>
#include <algorithm>

// Durum tutabilen bir fonksiyon nesnesi (functor)
class Multiplier {
    int coefficient_;
public:
    explicit Multiplier(int k) : coefficient_(k) {}
    int operator()(int x) const { return x * coefficient_; }   // cagrIlabilir
};

int main() {
    Multiplier triple(3);
    std::cout << triple(10) << "\n";     // 30 -- nesneyi fonksiyon gibi cagIr

    std::vector<int> v = {1, 2, 3, 4};
    // Functor'u STL algoritmasIna gec
    std::transform(v.begin(), v.end(), v.begin(), Multiplier(10));
    for (int x : v) std::cout << x << " ";  // 10 20 30 40
    std::cout << "\n";

    // Lambda: aslInda derleyicinin urettigi isimsiz bir functor'dur
    auto add = [addend = 100](int x) { return x + addend; };
    std::cout << add(5) << "\n";           // 105
    return 0;
}
\end{lstlisting}

\bilgi{Bir lambda ifadesi, aslında derleyicinin arka planda ürettiği,
\code{operator()} aşırı yüklenmiş isimsiz bir sınıftır. Yakalama listesi
(\code{[topla = 100]}) o sınıfın üye değişkenlerine dönüşür. Bu yüzden
functor'ları anlamak, lambdaların nasıl çalıştığını anlamaktır.}

\gsub{Dönüşüm Operatörü}

Bir sınıfı başka bir tipe \emph{örtük veya açık} dönüştürmek için dönüşüm
operatörü tanımlanır. \code{explicit} ile açık dönüşüm zorunlu kılınabilir.

\begin{lstlisting}
#include <iostream>

class Fraction {
    int numerator_, denominator_;
public:
    Fraction(int p, int q) : numerator_(p), denominator_(q) {}

    // double'a ORTUK donusum operatoru
    operator double() const {
        return static_cast<double>(numerator_) / denominator_;
    }
    // Bir sInIfIn 'bool baglamI'nda kullanIlmasI icin explicit tercih edilir
    explicit operator bool() const { return denominator_ != 0; }
};

int main() {
    Fraction half(1, 2);
    double d = half;                       // ortuk donusum -> 0.5
    std::cout << "double: " << d << "\n";
    std::cout << "1.5 + yarim = " << (1.5 + half) << "\n";  // 2.0

    if (half) std::cout << "gecerli kesir\n";   // explicit operator bool
    return 0;
}
\end{lstlisting}

\uyari{Örtük dönüşüm operatörleri (\code{operator double()} gibi) güçlüdür
ama tehlikelidir: beklenmedik yerlerde sessizce devreye girip belirsizliğe
ve hatalara yol açabilirler. Özellikle \code{operator bool()}'u neredeyse
her zaman \code{explicit} yapın; aksi halde sınıfınız kazara tam sayı
aritmetiğine veya karşılaştırmalara karışabilir. Genel tavsiye: dönüşüm
operatörlerini idareli ve mümkünse \code{explicit} kullanın.}

% #####################################################################
\gpart{12}{Modern C++ (C++17 / C++20)}
% #####################################################################

% =====================================================================
\gsec{Yapısal Bağlama, if-init ve Yararlı Tipler}

Modern C++, günlük kodu belirgin biçimde kısaltan ve güvenlileştiren pek
çok özellik getirdi. Bunlardan bazılarını (\code{auto}, range-for) zaten
kullandık; burada en pratik olanları toplu görelim.

\begin{lstlisting}
#include <iostream>
#include <map>
#include <tuple>
#include <optional>
#include <string>
#include <string_view>

// std::optional: "deger olmayabilir" durumunu tip guvenli ifade et
std::optional<int> parsePositive(int x) {
    if (x > 0) return x;      // deger var
    return std::nullopt;      // deger yok
}

// string_view: kopyalamadan string'e salt-okunur bakis (ucuz)
void printLength(std::string_view sv) {   // ne string kopyasi ne const char*
    std::cout << "\"" << sv << "\" -> " << sv.size() << " karakter\n";
}

int main() {
    // YAPISAL BAGLAMA (C++17): pair/tuple/struct'i tek satirda parcala
    std::map<std::string, int> scores{{"Ada", 90}, {"Alan", 85}};
    for (const auto& [name, score] : scores)     // .first/.second yerine
        std::cout << name << " = " << score << "\n";

    std::tuple<int, double, std::string> record{1, 3.14, "pi"};
    auto [id, value, label] = record;            // ucunu birden ac
    std::cout << id << " " << value << " " << label << "\n";

    // IF-INIT (C++17): kosul icinde degisken tanimla (kapsami dar tut)
    if (auto result = parsePositive(42); result.has_value())
        std::cout << "Pozitif: " << *result << "\n";

    if (auto result = parsePositive(-5); !result)
        std::cout << "Deger yok, varsayilan: " << result.value_or(0) << "\n";

    // string_view: gecici string, string literal, hepsi kopyasiz calisir
    std::string s = "merhaba dunya";
    printLength(s);
    printLength("sabit metin");
    return 0;
}
\end{lstlisting}

\ipucu{\code{std::string\_view}, bir string'i \emph{kopyalamadan} okumak
için idealdir; fonksiyon parametrelerinde \code{const std::string\&} yerine
sıkça tercih edilir çünkü string literalleri ve alt-dizeleri de kopyasız
kabul eder. Ancak dikkat: \code{string\_view} veriye \emph{sahip değildir};
gösterdiği string'ten daha uzun yaşarsa sarkan görünüm oluşur. Sahiplik
gerekiyorsa \code{std::string} kullanın.}

% =====================================================================
\gsec{Concepts: Şablonlara Kısıt (C++20)}

C++20 \emph{concepts}, şablon parametrelerine \emph{kısıtlar} koyar: ``bu
şablon yalnızca sayısal tiplerle çalışır'' gibi. Bu, hem niyeti belgeler
hem de yanlış tip kullanıldığında okunması imkânsız şablon hataları yerine
\emph{net} hata mesajları verir.

\begin{lstlisting}
#include <iostream>
#include <concepts>

// Concept tanimi: T aritmetik (sayisal) bir tip olmali
template <typename T>
concept Numeric = std::integral<T> || std::floating_point<T>;

// Yalnizca Numeric tipleri kabul eden sablon
template <Numeric T>
T square(T x) {
    return x * x;
}

// Alternatif kisit sozdizimi: requires ile
template <typename T>
    requires std::floating_point<T>
T average(T a, T b) {
    return (a + b) / 2;
}

int main() {
    std::cout << square(5) << "\n";        // int -> 25
    std::cout << square(2.5) << "\n";      // double -> 6.25
    // square("metin");  // DERLEME HATASI: net mesaj -> "Numeric saglanmadi"

    std::cout << average(3.0, 5.0) << "\n";  // 4.0
    return 0;
}
\end{lstlisting}

\begin{lstlisting}[style=shstyle]
g++ -std=c++20 concepts.cpp -o concepts
\end{lstlisting}

\bilgi{Concepts, C++ şablon programlamasının en büyük kalite sıçramalarından
biridir. \code{Numeric}, \code{std::integral}, \code{std::floating\_point},
\code{std::sortable} gibi kısıtlar sayesinde derleyici, bir şablonu yanlış
tiple kullandığınızda \emph{tam olarak hangi kısıtın sağlanmadığını}
söyler --- eski günlerin yüzlerce satırlık anlaşılmaz şablon hatalarına
son verir.}

% =====================================================================
\gsec{Ranges: Bileştirilebilir Algoritmalar (C++20)}

C++20 \emph{ranges} kütüphanesi, STL algoritmalarını hem daha kısa
(\code{begin()/end()} yazmadan) hem de \emph{bileştirilebilir} (composable)
hale getirir. ``Görünümler'' (views) sayesinde filtreleme, dönüştürme gibi
işlemleri bir boru hattı (pipeline) gibi zincirleyebilirsiniz --- üstelik
tembel (lazy) değerlendirmeyle.

\begin{lstlisting}
#include <iostream>
#include <vector>
#include <ranges>
#include <algorithm>

int main() {
    std::vector<int> numbers{1, 2, 3, 4, 5, 6, 7, 8, 9, 10};

    // Klasik algoritma: begin()/end() gerekmez
    std::ranges::sort(numbers);

    // Boru hatti: cift sayilari al -> karesini al (tembel degerlendirme)
    auto pipeline = numbers
        | std::views::filter([](int x){ return x % 2 == 0; })   // 2,4,6,8,10
        | std::views::transform([](int x){ return x * x; });     // 4,16,36,64,100

    std::cout << "Cift karelar: ";
    for (int x : pipeline) std::cout << x << " ";
    std::cout << "\n";

    // views::take, views::reverse gibi baska gorunumler
    std::cout << "Ilk 3 (tersten): ";
    for (int x : numbers | std::views::reverse | std::views::take(3))
        std::cout << x << " ";              // 10 9 8
    std::cout << "\n";
    return 0;
}
\end{lstlisting}

\begin{lstlisting}[style=shstyle]
g++ -std=c++20 ranges.cpp -o ranges
\end{lstlisting}

\ipucu{Ranges'in \code{|} boru hattı sözdizimi, veri dönüşümlerini son
derece okunur kılar: ``filtrele, sonra dönüştür, sonra ilk 3'ünü al''
niyetini doğrudan ifade eder. Görünümler \emph{tembeldir}: ara diziler
oluşturmaz, elemanlar yalnızca gezinirken hesaplanır --- bu hem bellek hem
hız açısından verimlidir. Modern C++ kodunda döngü ve ara vektör
yığınlarının yerini giderek ranges alıyor.}

% =====================================================================
\gsec{Kapanış ve Altın Kurallar}

Bu rehber, C++'ı temellerinden modern özelliklerine kadar kapsadı:
sözdizimi, nesne yönelimi, bellek ve kaynak yönetimi, şablonlar, STL,
lambda, istisnalar ve C++20 yenilikleri. Aşağıda, özellikle kaynak yönetimi
ve sınıf tasarımı için özet altın kurallar --- günlük kodda en çok fark
yaratanlar:

\begin{center}
\begin{tabular}{@{}c>{\raggedright\arraybackslash}p{11.0cm}@{}}
\toprule
\# & \textbf{Altın Kural} \\
\midrule
1 & \textbf{Sıfır Kuralını} hedefleyin: kaynağı \code{unique\_ptr}/konteynerlere
sarın, özel üye fonksiyonları hiç yazmayın. \\
2 & Çıplak \code{new}/\code{delete} kullanmayın; \code{make\_unique}/\code{make\_shared} tercih edin. \\
3 & \textbf{Varsayılan olarak \code{unique\_ptr}}; sahiplik gerçekten
paylaşılmalıysa \code{shared\_ptr}; döngüyü kırmak için \code{weak\_ptr}. \\
4 & Beşin Kuralını uygularsanız \emph{beşini de} bilinçli tanımlayın; taşımaları
\code{noexcept} yapın. \\
5 & Kopya atama için \textbf{copy-and-swap} deyimini kullanın. \\
6 & Tek argümanlı yapıcıları ve \code{operator bool}'u \code{explicit} yapın. \\
7 & Durumu değiştirmeyen üye fonksiyonları \code{const} yapın (const-correctness). \\
8 & \code{nullptr} kullanın (\code{NULL}/\code{0} değil); sildikten sonra işaretçiyi \code{nullptr}'a atayın. \\
9 & Operatörleri sezgisel anlamlarına sadık kalarak aşırı yükleyin;
\code{operator+}'yı \code{operator+=} cinsinden yazın. \\
10 & Derleyici uyarılarını açın: \code{-Wall -Wextra}; bellek hataları için
\code{-fsanitize=address}. \\
\bottomrule
\end{tabular}
\end{center}

\bilgi{\textbf{Bellek hatalarını yakalama:} Bu rehberdeki manuel kaynak
yöneten örnekleri test ederken AddressSanitizer'ı kullanın; sızıntıları,
sarkan işaretçileri ve ikili silmeleri anında yakalar:}

\begin{lstlisting}[style=shstyle]
# AddressSanitizer + tum uyarilar ile derleme (gelistirme sirasinda)
g++ -std=c++20 -Wall -Wextra -g -fsanitize=address,undefined \
    program.cpp -o program
./program        # bellek hatalari calisirken raporlanIr

# Valgrind ile sizinti kontrolu (alternatif)
g++ -std=c++20 -g program.cpp -o program
valgrind --leak-check=full ./program
\end{lstlisting}

\vspace{1cm}
\begin{center}
{\color{ortaGri}\rule{0.5\linewidth}{0.8pt}}\\[6pt]
{\color{anaMavi}\large\bfseries Rehberin Sonu}\\[4pt]
{\color{ortaGri}\small C++17/20 \textbullet\ RAII \textbullet\ Beşin \& Sıfırın Kuralı \textbullet\ İyi kodlamalar!}
\end{center}

\end{document}
